% ==== IX ====
\yoschapter{IX}{Analytical Theory of Semi-groups}
\label{chap:IX}

The analytical theory of semi-groups of bounded linear operators in a $B$-space deals with the exponential functions in infinite dimensional function spaces. It is concerned with the problem of determining the most general bounded linear operator valued function $T(t)$, $t \geqq 0$, which satisfies the equations 
\[
T (t + s) = T (t) \cdot T (s), T (0) = I.
\]

The problem was investigated by E. \textsc{Hille} \cite{hille98} and K. \textsc{Yosida} \cite{yosida348} independently of each other around $1948$. They introduced the notion of the infinitesimal generator $A$ of $T(t)$ defined by 
\[
A = \slim _{t \downarrow 0} t ^{-1} (T (t) - I),
\]
and discussed the generation of $T(t)$ in terms of $A$ and obtained a characterization of the infinitesimal generator $A$ in terms of the spectral property of $A$. 

The basic result of the semi-group theory may be considered as a natural generalization of the theorem of M. H. \textsc{Stone} \cite{stone312} on one-parameter group of unitary operators in a Hilbert space, which will be explained in a later section. Applications of the theory to stochastic processes and to the integration of the equations of evolution, which include diffusion equations, wave equations and Schrödinger equations, will be discussed in Chapter~\ref{chap:XIV}.\yosstarfootnote{See also Supplementary Notes, p.~\pageref{notes:IX-XIV}.} 

In this chapter, we shall develop the theory of semi-groups of continuous linear operators in locally convex linear topological spaces rather than in Banach spaces. 


\yossection{1}{The Semi-group of Class \texorpdfstring{$(C_0)$}{(C0)}}
\label{sec:IX.1}

\begin{propositionn}[E. \textsc{Hille}] Let $X$ be a $B$-space, and $T_{t}, t \geqq 0$, a one-parameter family of bounded linear operators $\in L(X, X)$ satisfying the semi-group property
\begin{equation}\label{eq:IX.1.1}
T _{t} T _{s} = T _{t+s} \quad \text{ for } \quad t, s > 0.
\end{equation}

If $p(t) = \log \|T_t\|$ is bounded from above on the interval $(0, a)$ for each positive $a$, then 
\begin{equation}\label{eq:IX.1.2}
\lim\limits _{t \rightarrow \infty} t ^{-1} \log \| T _{t} \| = \inf\limits _{t>0} t ^{-1} \log \| T _{t} \|.
\end{equation}

\end{propositionn}
\begin{proof}
 We have $p(t + s) \leqq p(t) + p(s)$ from $\| T_{t+s} \| = \| T_t T_s \| \leqq \| T_t \| \cdot \| T_s \|$. Let $\beta = \inf\limits_{t>0} t^{-1} \cdot p(t)$. $\beta$ is either finite or $-\infty$. Suppose that $\beta$ is finite. We choose, for any $\varepsilon > 0$, a number $a > 0$ in such a way that $p(a) \leqq (\beta + \varepsilon)a$. Let $t > a$ and $n$ be an integer such that $na \leqq t < (n + 1)a$. Then 
\[
\begin{aligned}
&\beta \leqq \frac {p (t)}{t} \leqq \frac {p (n a)}{t} + \frac {p (t - n a)}{t} \leqq \frac {n a}{t} \frac {p (a)}{a} + \frac {p (t - n a)}{t} \leqq \frac {n a}{t} (\beta + \varepsilon) + \frac {p (t - n a)}{t}.
\end{aligned}
\]

By hypothesis, $p(t - na)$ is bounded from above as $t \to \infty$. Thus, letting $t \to \infty$ in the above inequality, we obtain $\lim\limits_{t \to \infty} t^{-1} p(t) = \beta$. The case $\beta = -\infty$ may be treated similarly. 
\end{proof}

\begin{definition}\label{def:IX.1.Definition.1} If $\{T_t; t \geqq 0\} \subseteq L(X, X)$ satisfy the conditions 
\begin{align}
T _{t} T _{s}& = T _{t+s} (\text{ for } t, s \geqq 0), \tag{1'}\label{eq:IX.1.1'} \\
T _{0}& = I, \label{eq:IX.1.3} \\
\slim _ {t \to t _{0}} T _{t} x& = T _ {t _{0}} x \text{ for each } t _{0} \geqq 0 \text{ and each } x \in X, \label{eq:IX.1.4}
\end{align}
then $\{T_t\}$ is called a semi-group of class $(C_0)$. 
\end{definition}

In virtue of the Proposition, we see that a semi-group $\{T_t\}$ of class $(C_0)$ satisfies the condition 
\begin{equation}\label{eq:IX.1.5}
\left\| T _{t} \right\| \leqq M e ^{\beta t} \quad (\text{ for } 0 \leqq t < \infty),
\end{equation}
with constants $M > 0$ and $\beta < \infty$. 

The proof is easy. We have only to show that $\|T_t\|$ is, for any interval $(0, a)$ with $\infty > a > 0$, bounded on $(0, a)$. Assume the contrary and let there exist a sequence $\{t_n\} \subseteq (0, a)$ such that $\|T_{t_{n}}\| > n$ and $\lim\limits_{n \to \infty} t_n = t_0 \leqq a$. By the resonance theorem, $\|T_{t_{n}}x\|$ must be unbounded at least for one $x \in X$, which surely contradicts the strong continuity condition \eqref{eq:IX.1.4}. 

\begin{remarku} By multiplying $e^{-\beta t}$, we may assume that a semi-group $\{T_{t}\}$ of class $(C_{0})$ is equi-bounded: 
\begin{equation}\label{eq:IX.1.6}
\left\| T _{t} \right\| \leqq M \quad (\text{ for } 0 \leqq t < \infty).
\end{equation}
\end{remarku}

If, in particular, $M$ is $\leqq 1$, i.e., if 
\begin{equation}\label{eq:IX.1.7}
\left\| T _{t} \right\| \leqq 1 \quad (\text{ for } 0 \leqq t < \infty),
\end{equation}
then the semi-group $\{T_t\}$ is called a contraction semi-group of class $(C_0)$. As for the strong continuity condition \eqref{eq:IX.1.4}, we have the following 

\begin{theoremu} Let a family $\{T_t; t \geqq 0\}$ of operators $\in L(X, X)$ satisfy \eqref{eq:IX.1.1} and \eqref{eq:IX.1.3}. Then condition \eqref{eq:IX.1.4} is equivalent to the condition 
\begin{equation}\label{eq:IX.1.8}
\wlim _{t \downarrow 0} T _{t} x = x \quad \text{ for every } \quad x \in X.
\end{equation}
\end{theoremu}

\begin{proof}
 Suppose that \eqref{eq:IX.1.8} is satisfied. Let $x_0$ be any fixed element of $X$. We shall show that $s\text{- }\lim\limits_{t \to t_0} T_t x_0 = T_{t_0} x_0$ for each $t_0 \geqq 0$. Consider the function $x(t) = T_t x_0$. For each $t_0 \geqq 0$, $x(t)$ is weakly right continuous at $t_0$, because $w\text{- }\lim\limits_{t \downarrow t_0} T_t x = w\text{- }\lim\limits_{h \downarrow 0} T_h T_{t_0} x_0 = T_{t_0} x_0$. We next prove that $\|T_t\|$ is bounded in a vicinity of $t = 0$. For, otherwise, there would exist a sequence $\{t_n\} \in$ such that $t_n \downarrow 0$ and $\lim\limits_{n \to \infty} \|T_{t_{n}} x_0\| = \infty$, contrary to the resonance theorem implied by the weak right continuity of $x(t) = T_t x_0$. Thus by \eqref{eq:IX.1.1}, we see that $T_t x_0 = x(t)$ is bounded on any compact interval of $t$. Moreover, $x(t)$ is weakly measurable. For, a right continuous real-valued function $f(t)$ is Lebesgue measurable, as may be proved from the fact that, for any $\alpha$, the set $\{t; f(t) < \alpha\}$ is representable as the union of intervals of positive length. Next let $\{t_n\}$ be the totality of positive rational numbers, and consider finite linear combinations $\sum\limits_j \beta_j x(t_j)$ where $\beta_j$ are rational numbers (if $X$ is a complex linear space, we take $\beta_j = a_j + ib_j$ with rational coefficients $a_j$ and $b_j$). These elements form a countable set $M = \{x_n\}$ such that $\{x(t); t \geqq 0\}$ is contained in the strong closure of $M$. If otherwise, there would exist a number $t'$ such that $x(t')$ does not belong to $M^a$. But, being a closed linear subspace of $X$, $M^a$ is weakly closed by Theorem~\ref{thm:V.1.Theorem.11} in Chapter~\ref{chap:V}, \ref{sec:V.1}; consequently, the condition $x(t') \in M^a$ is contradictory to the weak right continuity of $x(t)$, i.e., to $x(t') = w\text{- }\lim\limits_{t_n \downarrow t'} x(t_n)$. 

We may thus apply Pettis' theorem in Chapter~\ref{chap:V}, \ref{sec:V.4} to the effect that $x(t)$ is strongly measurable, and so by the boundedness of $\|x(t)\|$ on any compact interval of $t$, we may define the Bochner integral $\int_{\alpha}^{\beta} x(t) \dd t$ and we have $\left\| \int_{\alpha}^{\beta} x(t) \dd t \right\| \leqq \int_{\alpha}^{\beta} \|x(t)\| \dd t$ for $0 \leqq \alpha < \beta < \infty$. By virtue of the strong continuity in $s$ of the integral $\int_{\alpha}^{\beta} x(t + s) \dd t = \int_{\alpha + s}^{\beta + s} x(t) \dd t$, which is implied by the boundedness of $x(t)$ on any compact interval of $t$, N. \textsc{Dunford} \cite{dunford46} proved that $x(t)$ is strongly continuous in $t > 0$. We shall follow his proof. 

Let $0 \leqq \alpha < \eta < \beta < \xi - \varepsilon < \xi$ with $\varepsilon > 0$. Since $x(\xi) = T_{\xi}x_0 = T_{\eta}T_{\xi - \eta}x_0 = T_{\eta}x(\xi - \eta)$, we have 
\[
(\beta - \alpha) x (\xi) = \int_{\alpha} ^{\beta} x (\xi) \dd \eta = \int_{\alpha} ^{\beta} T _{\eta} x (\xi - \eta) \dd \eta ,
\]
and so, by $\sup\limits_{\alpha \leqq \eta \leqq \beta}\|T_{\eta}\| < \infty$ which is implied by \eqref{eq:IX.1.1} and \eqref{eq:IX.1.3}, combined with the boundedness of $\|T_t\|$ near $t = 0$, we obtain 
\begin{align}
(\beta - \alpha) \{x (\xi \pm \varepsilon) - x (\xi) \}& = \int_{\alpha} ^{\beta} T _{\eta} \{x (\xi \pm \varepsilon - \eta) - x (\xi - \eta) \} \dd \eta , \notag \\
(\beta - \alpha) \| x (\xi \pm \varepsilon) - x (\xi) \|& \leqq \sup\limits _{\alpha \leqq \eta \leqq \beta} \| T _{\eta} \| \cdot \int_{\xi - \beta} ^{\xi - \alpha} \| x (\tau \pm \varepsilon) - x (\tau) \| \dd \tau . \notag
\end{align}

The right hand side tends to zero as $\varepsilon \downarrow 0$, as may be seen by approximating $x(\tau)$ by finitely-valued functions. 

We have thus proved that $x(t)$ is strongly continuous in $t > 0$. To prove the strong continuity of $x(t)$ at $t = 0$, we proceed as follows. For any positive rational number $t_n$, we have $T_t x(t_n) = T_t T_{t_{n}} x_0 = T_{t+t_n} x_0 = x(t + t_n)$. Hence, by the strong continuity of $x(t)$ for $t > 0$ proved above, we have $\slim\limits_{t \downarrow 0} T_t x(t_n) = x(t_n)$. Since each $x_m \in M$ is a finite linear combination of $x(t_n)'s$, we have $\slim\limits_{t \downarrow 0} T_t x_m = x_m (m = 1, 2, \ldots)$. On the other hand, we have, for any $t \in [0, 1]$, 
\begin{align}
\left\| x (t) - x _{0} \right\|& \leqq \left\| T _{t} x _{m} - x _{m} \right\| + \left\| x _{m} - x _{0} \right\| + \left\| T _{t} \left(x _{0} - x _{m}\right) \right\| \notag \\
&\leqq \left\| T _{t} x _{m} - x _{m} \right\| + \left\| x _{m} - x _{0} \right\| + \sup\limits _{0<t<1} \left\| T _{t} \right\| \cdot \left\| x _{0} - x _{m} \right\|. \notag
\end{align}

Hence $\varlimsup\limits_{t\downarrow 0}\|x(t) - x_0\|\leqq (1 + \sup\limits_{0\leqq t\leqq 1}\|T_t\|)\|x_m - x_0\|,$ and so $\slim\limits_{t\downarrow 0}x(t)$ $= x_0$ by $\inf\limits_{x_m\in M}\|x_0 - x_m\| = 0.$ 
\end{proof}


\yossection[2. The Equi-continuous Semi-group of Class \texorpdfstring{$(C_0)$}{(C0)}]{2}{The Equi-continuous Semi-group of Class \texorpdfstring{$(C_0)$}{(C0)} in Locally Convex Spaces. Examples of Semi-groups}
\label{sec:IX.2}

Suggested by the preceding section, we shall pass to a more general class of semi-groups. 

\begin{definitionu} Let $X$ be a locally convex linear topological space, and $\{T_t; t \geqq 0\}$ be a one-parameter family of continuous linear operators $\in L(X, X)$ such that 
\begin{align}
T _{t} T _{s}& = T _{t+s}, T _{0} = I, \label{eq:IX.2.1} \\
\lim\limits _ {t \rightarrow t _{0}} T _{t} x& = T _ {t _{0}} x \quad \text{ for any } \quad t _{0} \geqq 0 \quad \text{ and } \quad x \in X, \label{eq:IX.2.2}
\end{align}
the family of mappings $\{T_t\}$ is equi-continuous in $t$, 
i.e., for any continuous semi-norm $\phi$ on $X$, there 
exists a continuous semi-norm $q$ on $X$ such that 
\begin{equation}\label{eq:IX.2.3}
p \left(T _{t} x\right) \leqq q (x) \text{ for all } t \geqq 0 \text{ and all } x \in X.
\end{equation}
\end{definitionu}

Such a family $\{T_t\}$ is called an equi-continuous semi-group of class $(C_0)$. 

The semi-groups $\{T_{t}\}$ satisfying conditions \eqref{eq:IX.2.1}, \eqref{eq:IX.2.3}, ($4$) and ($6$) of Section~\hyperref[sec:IX.1]{$1$} are example of such equi-continuous semi-groups of class $(C_0)$. We shall give concrete examples. 

\begin{example}\label{ex:IX.2.Example.1} Let $C[0, \infty]$ be the space of bounded uniformly continuous real-valued (or complex-valued) functions on the interval $[0, \infty)$, and define $T_t, t \geqq 0$, on $C[0, \infty]$ into $C[0, \infty]$ by 
\[
(T _{t} x) (s) = x (t + s).
\]
\end{example}

Condition \eqref{eq:IX.2.1} is trivially satisfied. \eqref{eq:IX.2.2} follows from the uniform continuity of $x(s)$. Finally $\|T_t\| \leqq 1$ and so $\{T_t\}$ is a contraction semigroup of class $(C_0)$. In this example, we could replace $C[0,\infty]$ by $C[-\infty, \infty]$ or by $L^p(-\infty, \infty)$. 

\begin{example}\label{ex:IX.2.Example.2} Consider the space $C[-\infty, \infty]$. Let 
\[
N _{t} (u) = (2 \pi t) ^{-1/2} e ^ {- u ^{2} / 2 t}, - \infty < u < \infty , t > 0,
\]
which is the Gaussian probability density. Define $T_{t}, t \geqq 0$, on $C[-\infty, \infty]$ into $C[-\infty, \infty]$ by 
\[
\begin{aligned} (T _{t} x) (s) & = \int_{- \infty} ^{\infty} N _{t} (s - u) x (u) \dd u, \text{ for } t > 0, \\ & = x (s), \text{ for } t = 0. \end{aligned}
\]
\end{example}

Each $T_{t}$ is continuous, since, by $\int_{-\infty}^{\infty} N_{t}(s - u) \dd u = 1$, 
\[
\left\| T _{t} x \right\| \leqq \left\| x \right\| \int_{- \infty} ^{\infty} N _{t} (s - u) \dd u = \left\| x \right\|.
\]
$T_{0}=I$ by the definition, and the semi-group property $T_{t}T_{s}=T_{t+s}$ is a consequence from the well-known formula concerning the Gaussian probability distribution: 
\[
\frac {1}{\sqrt {2 \pi (t + t ^{\prime})}} e ^ {- u ^{2} / 2 (t + t ^{\prime})} = \frac {1}{\sqrt {2 \pi t}} \frac {1}{\sqrt {2 \pi t ^{\prime}}} \int_{- \infty} ^{\infty} e ^ {- (u - v) ^{2} / 2 t} e ^ {- v ^{2} / 2 t ^{\prime}} \dd v.
\]

This formula may be proved by applying the Fourier transformation on both sides, remembering ($10$) and \eqref{eq:VI.1.13} in Chapter~\ref{chap:VI}, \ref{sec:VI.1}. To prove the strong continuity in $t$ of $T_{t}$, we observe that $x(s) = \int_{-\infty}^{\infty} N_{t}(s - u) x(s) \dd u$. Thus 
\[
(T _{t} x) (s) - x (s) = \int_{- \infty} ^{\infty} N _{t} (s - u) (x (u) - x (s)) \dd u,
\]
which is, by the change of variable $(s - u) / \sqrt{t} = z$, equal to 
\[
\int_{- \infty} ^{\infty} N _{1} (z) \left(x (s - \sqrt {t} z) - x (s)\right) \dd z.
\]

By the uniform continuity of $x(s)$ on $(-\infty, \infty)$, there exists, for any $\varepsilon > 0$, a number $\delta = \delta(\varepsilon) > 0$ such that $|x(s_1) - x(s_2)| \leqq \varepsilon$ whenever $|s_1 - s_2| \leqq \delta$. Splitting the last integral, we obtain 
\[
\begin{aligned}
&\left| (T _{t} x) (s) - x (s) \right| \leqq \int_ {| \sqrt {t} z | \leqq \delta} N _{1} (z) \left| x (s - \sqrt {t} \cdot z) - x (s) \right| \dd z + \int_ {| \sqrt {t} z | > \delta} (\dots) \dd z
\\
&\leqq \varepsilon \int_ {| \sqrt {t} z | \leqq \delta} N _{1} (z) \dd z + 2 \| x \| \int_ {| \sqrt {t} z | > \delta} N _{1} (z) \dd z \leqq \varepsilon + 2 \| x \| \int_ {| \sqrt {t} z | > \delta} N _{1} (z) \dd z.
\end{aligned}
\]

The second term on the right tends to $0$ as $t \to 0$, since the integral $\int_{-\infty}^{\infty} N_1(z) \dd z$ converges. Thus $\lim\limits_{t \downarrow 0} \sup\limits_s |(T_t x)(s) - x(s)| = 0$ and hence $\slim\limits_{t \downarrow 0} T_t x = x$; consequently, by the Theorem in the preceding section, we have proved \eqref{eq:IX.2.2}. 

In the above example, we can replace $C[-\infty, \infty]$ by $L^p(-\infty, \infty)$. Consider, for example, $L^1(-\infty, \infty)$. In this case, we have 
\[
\left\| T _{t} x \right\| \leqq \iint_{- \infty} ^{\infty} N _{t} (s - u) | x (u) | \dd s \dd u \leqq \left\| x \right\|,
\]
by Fubini's theorem. As for the strong continuity, we have, as above, 
\[
\begin{aligned}
&\| T _{t} x - x \|  = \int_{- \infty} ^{\infty} \int N _{1} (z) (x (s - \sqrt {t} \cdot z) - x (s)) \dd z | \dd s \\
&\leqq \int_{- \infty} ^{\infty} N _{1} (z) \left[ \int_{- \infty} ^{\infty} | x (s - \sqrt {t} \cdot z) - x (s) | \dd s \right] \dd z \leqq 2 \int_{- \infty} ^{\infty} N _{1} (z) \dd z \cdot \| x \|.
\end{aligned}
\]

Hence, by the Lebesgue-Fatou lemma, we obtain 
\[
\varlimsup\limits_{t \downarrow 0} \| T _{t} x - x \| \leqq \int_{- \infty} ^{\infty} N _{1} (z) \left(\varlimsup\limits_{t \downarrow 0} \int_{- \infty} ^{\infty} | x (s - \sqrt {t} \cdot z) - x (s) | \dd s\right) \dd z = 0,
\]
because of the continuity in the mean of the Lebesgue integral, which may be proved by approximating $x(s)$ by finitely-valued functions. 

\begin{example}\label{ex:IX.2.Example.3} Consider $C[-\infty, \infty]$. Let $\lambda > 0$, $\mu > 0$. Define $T_t$, $t \geqq 0$, on $C[-\infty, \infty]$ into $C[-\infty, \infty]$ by 
\[
\left(T _{t} x\right) (s) = e ^{- \lambda t} \sum\limits_{k=0} ^{\infty} \frac {(\lambda t) ^{k}}{k !} x (s - k \mu).
\]
$3$. The Infinitesimal Generator of an Equi-continuous Semi-group We have 
\[
\begin{aligned}
\left(T _{w} (T _{t} x)\right) (s) &= e ^{- \lambda w} \sum\limits_{m=0} ^{\infty} \frac {(\lambda w) ^{m}}{m !} \left[ e ^{- \lambda t} \sum\limits_{k=0} ^{\infty} \frac {(\lambda t) ^{k}}{k !} x (s - k \mu - m \mu) \right]
\\
&= e ^{- \lambda (w + t)} \sum\limits_{p=0} ^{\infty} \frac {1}{p !} \left[ p! \sum\limits_{m=0} ^{p} \frac {(\lambda w) ^{m}}{m !} \frac {(\lambda t) ^{p-m}}{(p - m) !} x (s - p \mu) \right]
\\
&= e ^{- \lambda (w + t)} \sum\limits_{p=0} ^{\infty} \frac {1}{p !} (\lambda w + \lambda t) ^{p} x (s - p \mu) = (T _{w+t} x) (s).
\end{aligned}
\]
\end{example}

Thus it is easy to verify that $T_{t}$ is a contraction semi-group of class $(C_0)$. 


\yossection{3}{The Infinitesimal Generator of an Equi-continuous Semi-group of Class \texorpdfstring{$(C_0)$}{(C0)}}
\label{sec:IX.3}

Let $\{T_t; t \geqq 0\}$ be an equi-continuous semi-group of class $(C_0)$ defined on a locally convex linear topological space $X$ which we assume to be sequentially complete. We define the infinitesimal generator $A$ of $T_t$ by 
\begin{equation}\label{eq:IX.3.1}
A x = \lim\limits _{h \downarrow 0} h ^{-1} (T _{h} - I) x,
\end{equation}
i.e., $A$ is the linear operator whose domain is the set $D(A) = \{x \in X; \lim\limits_{h \downarrow 0} h^{-1}(T_h - I) x \text{ exists in } X\}$, and, for $x \in D(A)$, $Ax = \lim\limits_{h \downarrow 0} h^{-1}(T_h - I) x$. $D(A)$ is non-empty; it contains at least the vector $0$. Actually $D(A)$ is larger. We can prove 

\begin{theorem}\label{thm:IX.3.Theorem.1} $D(A)$ is dense in $X$. 
\end{theorem}

\begin{proof}
 Let $\varphi_{n}(s) = ne^{-ns}$, $n > 0$. Consider the linear operator $C_{\varphi_n}$ which is the Laplace transform of $T_{t}$ multiplied by $n$: 
\begin{equation}\label{eq:IX.3.2}
C _ {\varphi_{n}} x = \int_{0} ^{\infty} \varphi_{n} (s) T _{s} x \dd s \quad \text{ for } \quad x \in X,
\end{equation}
the integral being defined in the sense of Riemann. The ordinary procedure of defining the Riemann integral of numerical functions can be extended to a function with values in a locally convex, sequentially complete space $X$, using continuous semi-norms $p$ on $X$ in place of the absolute value of a number. The convergence of the improper integral is a consequence of the equi-continuity of $T_{t}$, the inequality 
\[
p \left(\varphi_{n} (s) T _{s} x\right) = n e ^{-ns} p \left(T _{s} x\right)
\]
and the sequential completeness of $X$. 


We see, by 
\[
\phi \left(C _ {\varphi_{n}} x\right) \leqq \int_{0} ^{\infty} n e ^{-ns} \phi \left(T _{s} x\right) \dd s \leqq \sup\limits _{s \geqq 0} \phi \left(T _{s} x\right),
\]
that the operator $C_{\varphi_n}$ is a continuous linear operator $\in L(X, X)$. We shall show that 
\begin{equation}\label{eq:IX.3.3}
R \left(C _ {\varphi_{n}}\right) \subseteq D (A) \quad \text{ for each } n > 0,
\end{equation}
and 
\begin{equation}\label{eq:IX.3.4}
\lim\limits _{n \rightarrow \infty} C _ {\varphi_{n}} x = x \quad \text{ for each } \quad x \in X.
\end{equation}

Then $\bigcup\limits_{n=1}^{\infty} R(C_{\varphi_n})$ will be dense in $X$, and a fortiori $D(A)$ will be dense in $X$. To prove \eqref{eq:IX.3.3}, we start with the formula 
\[
h ^{-1} (T _{h} - I) C _ {\varphi_{n}} x = h ^{-1} \int_{0} ^{\infty} \varphi_{n} (s) T _{h} T _{s} x \dd s - h ^{-1} \int_{0} ^{\infty} \varphi_{n} (s) T _{s} x \dd s.
\]

The change of the order: $T_{h}\int\limits_{0}^{\infty}\cdots=\int\limits_{0}^{\infty}T_{h}\cdots$ is justified, using the linearity and the continuity of $T_{h}$. Thus 
\[
\begin{aligned} h ^{-1} (T _{h} - I) C _ {\varphi_{n}} x & = h ^{-1} \int_{0} ^{\infty} \varphi_{n} (s) T _{s+h} x \dd s - h ^{-1} \int_{0} ^{\infty} \varphi_{n} (s) T _{s} x \dd s \\ & = \frac {e ^{nh} - 1}{h} n \int_{h} ^{\infty} e ^{- n \sigma} T _{\sigma} x \dd \sigma - \frac {1}{h} n \int_{0} ^{h} e ^{-ns} T _{s} x \dd s \\ & = \frac {e ^{nh} - 1}{h} \left\{C _ {\varphi_{n}} x - \int_{0} ^{h} n e ^{- n \sigma} T _{\sigma} x \dd \sigma \right\} - \frac {1}{h} \int_{0} ^{h} \varphi_{n} (s) T _{s} x \dd s. \end{aligned}
\]

By the continuity of $\varphi_{n}(s) T_{s}x$ in $s$, the second term on the right tends to $-\varphi_{n}(0) T_{0}x = -nx$ as $h \downarrow 0$. Similarly, the first term on the right tends to $nC_{\varphi_n}x$ as $h \downarrow 0$. Hence we have 
\begin{equation}\label{eq:IX.3.5}
A C _ {\varphi_{n}} x = n \left(C _ {\varphi_{n}} - I\right) x, x \in X.
\end{equation}

We next prove \eqref{eq:IX.3.4}. We have, by $\int_0^\infty ne^{-ns}\dd s = 1$, 
\begin{align}
C _ {\varphi_{n}} x - x& = n \int_{0} ^{\infty} e ^{-ns} (T _{s} x - x) \dd s, \notag \\
p \left(C _ {\varphi_{n}} x - x\right)& \leqq n \int_{0} ^{\infty} e ^{-ns} p \left(T _{s} x - x\right) \dd s = n \int_{0} ^{\delta} + n \int_{\delta} ^{\infty} = I _{1} + I _{2}, \text{ say }, \notag
\end{align}
where $\delta > 0$ is a positive number. For any $\varepsilon > 0$, we can choose, by the continuity in $s$ of $T_s x$, a $\delta > 0$ such that $p(T_s x - x) \leqq \varepsilon$ for $0 \leqq s \leqq \delta$. Then 
\[
I _{1} \leqq \varepsilon n \int_{0} ^{\delta} e ^{-ns} \dd s \leqq \varepsilon n \int_{0} ^{\infty} e ^{-ns} \dd s = \varepsilon .
\]

For a fixed $\delta > 0$, we have, by the equi-boundedness of $\{T_s x\}$ in $s \geqq 0$, 
\[
I _{2} \leqq n \int_{\delta} ^{\infty} e ^{-ns} (p (T _{s} x) + p (x)) \dd s \rightarrow 0 \quad \text{ as } \quad n \uparrow \infty .
\]

Hence we have proved \eqref{eq:IX.3.4}.
\end{proof}

\begin{definitionu} For $x \in X$, we define 
\begin{equation}\label{eq:IX.3.6}
D _{t} T _{t} x = \lim\limits _{h \rightarrow 0} h ^{-1} (T _{t+h} - T _{t}) x
\end{equation}
if the right hand limit exists. 
\end{definitionu}

\begin{theorem}\label{thm:IX.3.Theorem.2} If $x \in D(A)$, then $x \in D(D_t T_t)$ and 
\begin{equation}\label{eq:IX.3.7}
D _{t} T _{t} x = A T _{t} x = T _{t} A x, t \geqq 0.
\end{equation}
\end{theorem}

Thus, in particular, the operator $A$ is commutative with $T_{t}$. 

\begin{proof}
 If $x \in D(A)$, then we have, since $T_{t}$ is continuous linear, 
\[ \begin{aligned} T _{t} A x & = T _{t} \lim\limits _{h \downarrow 0} h ^{-1} (T _{h} - I) x = \lim\limits _{h \downarrow 0} h ^{-1} (T _{t} T _{h} - T _{t}) x = \lim\limits _{h \downarrow 0} h ^{-1} (T _{t+h} - T _{t}) x \\ & = \lim\limits _{h \downarrow 0} h ^{-1} (T _{h} - I) T _{t} x = A T _{t} x. \end{aligned}
\]

Thus, if $x \in D(A)$, then $T_{t}x \in D(A)$ and $T_{t}Ax = AT_{t}x = \lim\limits_{h\downarrow 0} h^{-1}(T_{t+h} - T_{t})x$. We have thus proved that the right derivative of $T_{t}x$ exists for each $x \in D(A)$. We shall show that the left derivative also exists and is equal to the right derivative. 

For this purpose, take any $f_0 \in X'$. Then, for a fixed $x \in D(A)$, the numerical function $f_0(T_t x) = \langle T_t x, f_0 \rangle$ is continuous in $t \geqq 0$ and has right derivative $d^+ f_0(T_t x) / \dd t$ which is equal to $f_0(A T_t x) = f_0(T_t A x)$ by what we have proved above. Hence $d^+ f_0(T_t x) / \dd t$ is continuous in $t$. We shall prove below a well-known Lemma: if one of the Dini derivatives 
\[
\overline {{{D}}} ^{+} f (t), \underline {{{D}}} ^{+} f (t), \overline {{{D}}} ^{-} f (t) \quad \text{ and } \quad \underline {{{D}}} ^{-} f (t)
\]
of a continuous real-valued function $f(t)$ is finite and continuous, then $f(t)$ is differentiable and the derivative is, of course, continuous and equals $\overline{D}^{\pm}f(t)$. Thus $f_0(T_tx)$ is differentiable in $t$ and 
\[
\begin{aligned} f _{0} (T _{t} x - x) & = f _{0} (T _{t} x) - f _{0} (T _{0} x) = \int_{0} ^{t} d ^{+} f _{0} (T _{s} x) / \dd s \cdot \dd s = \int_{0} ^{t} f _{0} (T _{s} A x) \dd s \\ & = f _{0} \left(\int_{0} ^{t} T _{s} A x \dd s\right). \end{aligned}
\]

Since $f_0 \in X'$ was arbitrary, we must have 
\[
T _{t} x - x = \int_{0} ^{t} T _{s} A x \dd s \quad \text{ for each } \quad x \in D (A).
\]

Since $T_{s}Ax$ is continuous in $s$, it follows that $T_{t}x$ is differentiable in $t$ in the topology of $X$ and 
\[
D _{t} T _{t} x = \lim\limits _{h \rightarrow 0} h ^{-1} \int_{t} ^{t+h} T _{s} A x \dd s = T _{t} A x.
\]

Thus we have proved \eqref{eq:IX.3.7}. 
\end{proof}

\begin{proof}[Proof of the Lemma.]
We first prove that the condition: $\overline{D}^{+}f(t) \geqq 0$ for $a \leqq t \leqq b$ implies that $f(b) - f(a) \geqq 0$. Assume the contrary, and let $f(b) - f(a) < -\varepsilon (b - a)$ with some $\varepsilon > 0$. Then, for $g(t) = f(t) - f(a) + \varepsilon (t - a)$, we have $\overline{D}^{+}g(a) = \overline{D}^{+}f(a) + \varepsilon > 0$, and so, by $g(a) = 0$, we must have $g(t_0) > 0$ for some $t_0 > a$ near $a$. By the continuity of $g(t)$ and $g(b) < 0$, there must exist a $t_1$ with $a < t_0 < t_1 < b$ such that $g(t_1) = 0$ and $g(t) < 0$ for $t_1 < t < b$. We have then $\overline{D}^{+}g(t_1) \leqq 0$ which surely contradicts the fact $\overline{D}^{+}g(t_1) = \overline{D}^{+}f(t_1) + \varepsilon > 0$.

By applying a similar argument to $f(t) - \alpha t$ and to $\beta t - f(t)$, we prove the following: if one of the Dini derivatives $Df(t)$ satisfies 
\[
\alpha \leqq D f (t) \leqq \beta \text{ on any interval } [ t _{1}, t _{2} ],
\]
then $\alpha \leqq (f(t_2) - f(t_1)) / (t_2 - t_1) \leqq \beta$. Hence, the suprema (and the infima) on $[t_1, t_2]$ of the four Dini derivatives of a continuous real-valued function $f(t)$ are the same. Thus, in particular, if one of the Dini derivatives of a continuous real-valued function $f(t)$ is continuous on $[t_1, t_2]$, then the four Dini derivatives of $f(t)$ must coincide on $[t_1, t_2]$. 
\end{proof}


\yossection{4}{The Resolvent of the Infinitesimal Generator \texorpdfstring{$A$}{A}}
\label{sec:IX.4}

\begin{theorem}\label{thm:IX.4.Theorem.1} If $n > 0$, then the operator $(nI - A)$ admits an inverse $R(n;A) = (nI - A)^{-1} \in L(X,X)$, and 
\begin{equation}\label{eq:IX.4.1}
R (n; A) x = \int_{0} ^{\infty} e ^{-ns} T _{s} x \dd s \quad \text{ for } \quad x \in X.
\end{equation}
\end{theorem}

In other words, positive real numbers belong to the resolvent set $\varrho(A)$ of $A$. 

\begin{proof}
 We first show that $(nI-A)^{-1}$ exists. Suppose that there exists an $x_{0}\neq0$ such that $(nI-A)x_{0}=0$, that is, $Ax_{0}=nx_{0}$. Let $f_{0}$ be a continuous linear functional $\in X'$ such that $f_{0}(x_{0})=1$, and set $\varphi(t)=f_{0}(T_{t}x_{0})$. Since $x_{0}\in D(A)$, $\varphi(t)$ is differentiable by Theorem~\ref{thm:IX.4.Theorem.2} in the preceding section and 
\[
\dd \varphi (t) / \dd t = f _{0} \left(D _{t} T _{t} x _{0}\right) = f _{0} \left(T _{t} A x _{0}\right) = f _{0} \left(T _{t} n x _{0}\right) = n \varphi (t). \]

If we solve this differential equation under the initial condition $\varphi(0)=f_{0}(x_{0})=1$, we get $\varphi(t)=e^{nt}$. But, $\varphi(t)=f_{0}(T_{t}x_{0})$ is bounded in $t$, because of the equi-boundedness of $T_{t}x_{0}$ in $t\geqq0$ and the continuity of the functional $f_{0}$. This is a contradiction, and so the inverse $(nI-A)^{-1}$ must exist. 

Since, by \eqref{eq:IX.4.5} of the preceding section, $A C_{\varphi_n}x = n(C_{\varphi_n} - I)x$, we have $(nI - A)C_{\varphi_n}x = nx$ for all $x \in X$. Thus $(nI - A)$ maps $R(C_{\varphi_n}) \subseteq D(A)$ onto $X$ in a one-one way. Hence, a fortiori, $(nI - A)$ must map $D(A)$ onto $X$ in a one-one way, in virtue of the existence of $(nI - A)^{-1}$. Therefore, we must have $R(C_{\varphi_n}) = D(A)$ and $(nI - A)^{-1} = n^{-1}C_{\varphi_n}$. 
\end{proof}
\begin{corollary}\label{cor:IX.4.Corollary.1} The right half plane of the complex $\lambda$-plane is in the resolvent set $\varrho(A)$ of $A$, and we have 
\begin{equation}\label{eq:IX.4.2}
R (\lambda ; A) x = (\lambda I - A) ^{-1} x = \int_{0} ^{\infty} e ^{- \lambda t} T _{t} x \dd t \text{ for } \operatorname{Re} (\lambda) > 0 \text{ and } x \in X.
\end{equation}
\end{corollary}

\begin{proof}
 For real fixed $\tau, \{e^{-i\tau t} T_t; t \geqq 0\}$, constitutes an equi-continuous semi-group of class $(C_0)$. It is easy to see that the infinitesimal generator of this semi-group is equal to $(A - i\tau I)$. Thus, for any $\sigma > 0$, the resolvent $R(\sigma + i\tau; A) = ((\sigma + i\tau)I - A)^{-1}$ exists and 
\begin{equation}\tag{2'}\label{eq:IX.4.2'}
R \left((\sigma + i \tau) I - A\right) x = \int_{0} ^{\infty} e ^{-(\sigma + i \tau) s} T _{s} x \dd s \quad \text{ for all } \quad x \in X. \end{equation}
\end{proof}

\begin{corollary}\label{cor:IX.4.Corollary.2} 
\begin{align}
D (A)& = R ((\lambda I - A) ^{-1}) = R (R (\lambda ; A)) \quad \text{ when } \quad \operatorname{Re} (\lambda) > 0, \label{eq:IX.4.3} \\
A R (\lambda ; A) x& = R (\lambda ; A) A x = (\lambda R (\lambda ; A) - I) x \text{ for } x \in D (A) \label{eq:IX.4.4} \\
A R (\lambda ; A) x& = (\lambda R (\lambda ; A) - I) x \quad \text{ for } \quad x \in X, \label{eq:IX.4.5} \\
\lim\limits _{\lambda \uparrow \infty} \lambda R (\lambda ; A) x& = x \quad \text{ for } \quad x \in X. \label{eq:IX.4.6}
\end{align}
\end{corollary}

\begin{proof}
 Clear from $R(\lambda; A) = (\lambda I - A)^{-1}$ and \eqref{eq:IX.4.4} of the preceding section. 
\end{proof}

\begin{corollary}\label{cor:IX.4.Corollary.3} The infinitesimal generator $A$ is a closed operator in the following sense (Cf. Chapter~\ref{chap:II}, \ref{sec:II.6}): 
\begin{align}
\text{ if } x _{h}& \in D (A) \quad \text{ and } \quad \lim\limits _{h \to \infty} x _{h} = x \in X, \lim\limits _{h \to \infty} A x _{h} = y \in X, \text{ then } \notag \\
x& \in D (A) \quad \text{ and } \quad A x = y. \notag
\end{align}
\end{corollary}

\begin{proof}
 Put $(I - A)x_{h} = z_{h}$. Then $\lim\limits_{h\to \infty}z_h = x - y$ and so, by the continuity of $(I - A)^{-1}$, $\lim\limits_{h\to \infty}x_h = \lim\limits_{h\to \infty}(I - A)^{-1}z_h = (I - A)^{-1}(x - y)$ that is, $x = (I - A)^{-1}(x - y)$, $(I - A)x = x - y$. This proves that $y = Ax$. 
\end{proof}

\begin{theorem}\label{thm:IX.4.Theorem.2} The family of operators 
\begin{equation}\label{eq:IX.4.7}
\{(\lambda R (\lambda ; A)) ^{n} \}
\end{equation}
is equi-continuous in $\lambda > 0$ and in $n = 0, 1, 2, \ldots$. 
\end{theorem}

\begin{proof}
 From the resolvent equation (Chapter~\ref{chap:VIII}, \ref{sec:VIII.2}, \eqref{eq:VIII.2.2}) 
\[
R (\mu ; A) - R (\lambda ; A) = (\lambda - \mu) R (\mu ; A) R (\lambda ; A),
\]
we obtain 
\[
\lim\limits _{\mu \rightarrow \lambda} (\mu - \lambda) ^{-1} (R (\mu ; A) - R (\lambda ; a)) x = \dd R (\lambda ; A) x / \dd \lambda = - R (\lambda ; A) ^{2} x, x \in X. \]

To derive the above formula, we have to appeal to \eqref{eq:IX.4.2} in order to show that $\lim\limits_{\mu \to \lambda} R(\mu; A)y = R(\lambda; A)y, y \in X$. 


Therefore, $R(\lambda; A) x$ is infinitely differentiable with respect to $\lambda$ when $\operatorname{Re}(\lambda) > 0$ and 
\begin{equation}\label{eq:IX.4.8}
d ^{n} R (\lambda ; A) x / \dd \lambda^{n} = (- 1) ^{n} n! R (\lambda ; A) ^{n+1} x \quad (n = 0, 1, 2, \dots).
\end{equation}

On the other hand, we have, by differentiating \eqref{eq:IX.4.2} $n$-times with respect to $\lambda$, 
\begin{equation}\label{eq:IX.4.9}
d ^{n} R (\lambda ; A) x / \dd \lambda^{n} = \int_{0} ^{\infty} e ^{- \lambda t} (- t) ^{n} T _{t} x \dd t.
\end{equation}

Here the differentiation under the integral sign is justified since $\{T_t x\}$ is equi-bounded in $t$ and $\int_0^\infty e^{-\lambda t} t^n \dd t = (n!) / \lambda^{n+1}$ when $\operatorname{Re}(\lambda) > 0$. Hence 
\begin{equation}\label{eq:IX.4.10}
(\lambda R (\lambda ; A)) ^{n+1} x = \frac {\lambda^{n+1}}{n !} \int_{0} ^{\infty} e ^{- \lambda t} t ^{n} T _{t} x \dd t \text{ for } x \in X \text{ and } \operatorname{Re} (\lambda) > 0,
\end{equation}
and so, for any continuous semi-norm $p$ on $X$ and $\lambda > 0$, $n > 0$, 
\begin{equation}\label{eq:IX.4.11}
\phi ((\lambda R (\lambda ; A)) ^{n+1} x) \leqq \frac {\lambda^{n+1}}{n !} \int_{0} ^{\infty} e ^{- \lambda t} t ^{n} \dd t \cdot \sup\limits _{t \geqq 0} \phi (T _{t} x) = \sup\limits _{t \geqq 0} \phi (T _{t} x).
\end{equation}

This proves Theorem~\ref{thm:IX.4.Theorem.2} by the equi-continuity of $\{T_t\}$ in $t$. 
\end{proof}


\yossection{5}{Examples of Infinitesimal Generators}
\label{sec:IX.5}

We first define, for $n > 0$, 
\begin{equation}\label{eq:IX.5.1}
J _{n} = (I - n ^{-1} A) ^{-1} = n R (n; A),
\end{equation}
so that 
\begin{equation}\label{eq:IX.5.2}
A J _{n} = n (J _{n} - I).
\end{equation}

\begin{example}\label{ex:IX.5.Example.1} $(T_{t}x)(s) = x(t + s)$ on $C[0,\infty]$. 
\end{example}

Writing $y_{n}(s) = (J_{n}x)(s) = n\int \limits_{0}^{\infty}e^{-nt}x(t + s)\dd t = n\int \limits_{s}^{\infty}e^{-n(t-s)}x(t)\dd t,$ 
we obtain $y_n'(s) = -n e^{-n(s-s)} x(s) + n^2 \int_s^{\infty} e^{-n(t-s)} x(t) \dd t = -n x(s) + n y_n(s)$. Comparing this with the general formula \eqref{eq:IX.5.2}: 
\[
(A J _{n} x) (s) = n ((J _{n} - I) x) (s),
\]
we obtain $A y_{n}(s) = y_{n}^{\prime}(s)$. Since $R(J_{n}) = R(R(n;A)) = D(A)$, we have $A y(s) = y^{\prime}(s)$ for every $y \in D(A)$. 

Conversely, let $y(s)$ and $y'(s)$ both belong to $C[0, \infty]$. We will then show that $y \in D(A)$ and $A y(s) = y'(s)$. To this purpose, define $x(s)$ by 
\[
y ^{\prime} (s) - n y (s) = - n x (s).
\]

Setting $(J_{n}x)(s) = y_{n}(s)$, we obtain, as shown above, 
\[
y _{n} ^{\prime} (s) - n y _{n} (s) = - n x (s).
\]

Hence $w(s) = y(s) - y_{n}(s)$ satisfies $w'(s) = nw(s)$ and so $w(s) = Ce^{ns}$. But, $w$ must belong to $C[0,\infty]$ and this is possible only if $C = 0$. Hence $y(s) = y_{n}(s) \in D(A)$ and $Ay(s) = y'(s)$. 

Therefore, the domain $D(A)$ of $A$ is precisely the set of functions $y \in C[0, \infty]$ whose first derivative also $\in C[0, \infty]$, and for such a function $y$ we have $A y = y'$. We have thus characterized the differential operator $\dd / \dd s$ as the infinitesimal generator of the semi-group associated with the operation of translation by $t$ on the function space $C[0, \infty]$. 

\begin{example}\label{ex:IX.5.Example.2} We shall give a characterization of the second derivative $d^{2}/\dd s^{2}$ as the infinitesimal generator of the semi-group associated with the integral operator by Gaussian kernel. The space is $C[-\infty,\infty]$ and 
\[
(T _{t} x) (s) = (2 \pi t) ^{-1/2} \int_{- \infty} ^{\infty} e ^ {- (s - v) ^{2} / 2 t} x (v) \dd v \text{ if } t > 0, = x (s) \text{ if } t = 0.
\]
\end{example}

We have 
\[
\begin{aligned}
&y _{n} (s) = (J _{n} x) (s) = \int_{- \infty} ^{\infty} x (v) \left[ \int_{0} ^{\infty} n (2 \pi t) ^{-1/2} e ^ {- n t - (s - v) ^{2} / 2 t} \dd t \right] \dd v
\\
&= \int_{- \infty} ^{\infty} x (v) \left[ 2 \sqrt {n} \int_{0} ^{\infty} (2 \pi) ^{-1/2} e ^ {- \sigma^{2} - n (s - v) ^{2} / 2 \sigma^{2}} \dd \sigma \right] \dd v (\text{ by substituting } t = \sigma^{2} / n).
\end{aligned}
\]

Assuming, for the moment, the formula 
\begin{equation}\label{eq:IX.5.3}
\int_{0} ^{\infty} e ^ {- (\sigma^{2} + c / \sigma^{2})} \dd \sigma = \frac {\sqrt {\pi}}{2} e ^{-2c}, c = \sqrt {n} | s - v | / \sqrt {2} > 0,
\end{equation}
we get 
\[
y _{n} (s) = \int_{- \infty} ^{\infty} x (v) (n / 2) ^{1/2} e ^ {- \sqrt {2 n} | s - v |} \dd v.
\]
$x(v)$ being continuous and bounded, we can differentiate twice and obtain 
\[
\begin{aligned}
y _{n} ^{\prime} (s) &= n \int_{s} ^{\infty} x (v) e ^ {- \sqrt {2 n} (v - s)} \dd v - n \int_{- \infty} ^{s} x (v) e ^ {- \sqrt {2 n} (s - v)} \dd v,
\\
y _{n} ^{\prime \prime} (s) &= n \left\{- x (s) - x (s) + \sqrt {2} \sqrt {n} \int_{s} ^{\infty} x (v) e ^ {- \sqrt {2 n} (s - v)} \dd v \right.
\\
&\quad \left. + \sqrt {2} \sqrt {n} \int_{- \infty} ^{s} x (v) e ^ {- \sqrt {2 n} (s - v)} \dd v \right\}
\\
&= - 2 n x (s) + 2 n y _{n} (s).
\end{aligned}
\]

Comparing this with the general formula \eqref{eq:IX.5.2}: 
\[
(A y _{n}) (s) = (A J _{n} x) (s) = n ((J _{n} - I) x) (s) = n (y _{n} (s) - x (s)),
\]
we find that $A y_{n}(s) = y_{n}^{\prime \prime}(s) / 2$. Since $R(J_{n}) = R(R(n;A)) = D(A)$, we  
have proved that 
\[
A y (s) = y ^{\prime \prime} (s) / 2 \quad \text{ for every } \quad y \in D (A).
\]

Conversely, let $y(s)$ and $y''(s)$ both belong to $C[-\infty, \infty]$. Define $x(s)$ by 
\[
y ^{\prime \prime} (s) - 2 n y (s) = - 2 n x (s).
\]

Setting $y_{n}(s) = (J_{n}x)(s)$, we obtain, as shown above, 
\[
y _{n} ^{\prime \prime} (s) - 2 n y _{n} (s) = - 2 n x (s).
\]

Hence $w(s) = y(s) - y_n(s)$ satisfies $w''(s) - 2nw(s) = 0$, and so $w(s) = C_1 e^{\sqrt{2n}s} + C_2 e^{-\sqrt{2n}s}$. This function cannot be bounded unless both $C_1$ and $C_2$ are zero. Thus $y(s) = y_n(s)$, and so $y \in D(A)$, $(Ay)(s) = y''(s)/2$. 

Therefore, the differential operator $\frac{1}{2}\frac{a^{2}}{\dd s^{2}}$ is characterized as the infinitesimal generator of the semi-group associated with the integral transform by the Gaussian kernel on the function space $C[-\infty,\infty]$. 

The proof of \eqref{eq:IX.5.3}. We start with the well-known formula 
\[
\int_{0} ^{\infty} e ^ {- x ^{2}} \dd x = \sqrt {\pi} / 2.
\]

Setting $x = \sigma - c / \sigma$, we obtain 
\[
\begin{aligned}
\frac {\sqrt {\pi}}{2} &= \int_ {\sqrt {c}} ^{\infty} e ^ {- (\sigma - c / \sigma) ^{2}} (1 + c / \sigma^{2}) \dd \sigma = e ^{2c} \int_ {\sqrt {c}} ^{\infty} e ^ {- (\sigma^{2} + c ^{2} / \sigma^{2})} (1 + c / \sigma^{2}) \dd \sigma
\\
&= e ^{2c} \left\{\int_ {\sqrt {c}} ^{\infty} e ^ {- (\sigma^{2} + c ^{2} / \sigma^{2})} \dd \sigma + \int_ {\sqrt {c}} ^{\infty} e ^ {- (\sigma^{2} + c ^{2} / \sigma^{2})} \frac {c}{\sigma^{2}} \dd \sigma \right\}.
\end{aligned}
\]

Setting $\sigma = c / t$ in the last integral, we obtain 
\[
\frac {\sqrt {\pi}}{2} = e ^{2c} \left\{\int_ {\sqrt {c}} ^{\infty} e ^ {- (\sigma^{2} + c ^{2} / \sigma^{2})} \dd \sigma - \int_ {\sqrt {c}} ^{0} e ^ {- (c ^{2} / t ^{2} + t ^{2})} \dd t \right\} = e ^{2c} \int_{0} ^{\infty} e ^ {- (t ^{2} + c ^{2} / t ^{2})} \dd t.
\]

Exercise. Show that the infinitesimal generator $A$ of the semi-group $\{T_t\}$ on $C[-\infty, \infty]$ given by 
\[
(T _{t} x) (s) = e ^{- \lambda t} \sum\limits_{k=0} ^{\infty} \frac {(\lambda t) ^{k}}{k !} x (s - k \mu) \quad (\lambda , \mu > 0),
\]
is the difference operator $A$: 
\[
(A x) (s) = \lambda \left\{x (s - \mu) - x (s) \right\}.
\]


\yossection{6}{The Exponential of a Continuous Linear Operator whose Powers are Equi-continuous}
\label{sec:IX.6}

\begin{propositionu} Let $X$ be a locally convex, sequentially complete, linear topological space. Let $B$ be a continuous linear operator $\in L(X, X)$ such that $\{B^k; k = 1, 2, \ldots\}$ is equi-continuous. Then, for each $x \in X$, the series 
\begin{equation}\label{eq:IX.6.1}
\sum\limits_{k=0} ^{\infty} (k!) ^{-1} (t B) ^{k} x (t \geqq 0)
\end{equation}
converges. 
\end{propositionu}

\begin{proof}
 For any continuous semi-norm $p$ on $X$, there exists, by the equi-continuity of $\{B^k\}$, a continuous semi-norm $q$ on $X$ such that $p(B^k x) \leqq q(x)$ for all $k \geqq 0$ and $x \in X$. Hence 
\[
\phi \left(\sum\limits_{k=n} ^{m} (t B) ^{k} x / k!\right) \leqq \sum\limits_{k=n} ^{m} t ^{k} \phi (B ^{k} x) / k! \leqq q (x) \cdot \sum\limits_{k=n} ^{m} t ^{k} / k!. \]

Therefore $\left\{\sum\limits_{k=0}^{n}(tB)^{k}x/k!\right\}$ is a Cauchy sequence in the sequentially complete space $X$. The limit of this sequence will be denoted by \eqref{eq:IX.6.1}. 
\end{proof}

\begin{corollary}\label{cor:IX.6.Corollary.1} The mapping $x \to \sum\limits_{k=0}^{\infty} (tB)^k x / k!$ defines a continuous linear operator which we shall denote by $\exp(tB)$. 
\end{corollary}

\begin{proof}
 By the equi-continuity of $\{B^k\}$, we can prove that $B_n = \sum\limits_{k=0}^{n} (tB)^k / k!$ ( $n = 0, 1, 2, \ldots$) are equi-continuous when $t$ ranges over any compact interval. In fact, we have 
\[
p (B _{n} x) \leqq \sum\limits_{k=0} ^{n} t ^{k} p (B ^{k} x) / k! \leqq q (x) \cdot \sum\limits_{k=0} ^{n} t ^{k} / k! \leqq e ^{t} \cdot q (x). \]

Hence the limit $\exp (tB)$ satisfies 
\begin{equation}\label{eq:IX.6.2}
p (\exp (t B) x) \leqq \exp (t) \cdot q (x) (t \geqq 0).
\end{equation}
\end{proof}

\begin{corollary}\label{cor:IX.6.Corollary.2} Let $B$ and $C$ be two continuous linear operators $\in L(X, X)$ such that $\{B^k\}$ and $\{C^k\}$ are equi-continuous. If, moreover, $BC = CB$, then we have 
\begin{equation}\label{eq:IX.6.3}
\exp (t B) \exp (t C) = \exp (t (B + C)).
\end{equation}
\end{corollary}

\begin{proof}
 We have 
\[
p \left((B + C) ^{k} x\right) \leqq \sum\limits_{s=0} ^{k} k C _{s} p \left(B ^{k-s} C ^{s} x\right) \leqq \sum\limits_{s=0} ^{k} k C _{s} q \left(C ^{s} x\right) \leqq 2 ^{k} \sup\limits _{0 \leqq s} q \left(C ^{s} x\right) \]

Hence $\{2^{-k}(B + C)^k\}$ is equi-continuous, and we can define $\exp(t(B + C))$. By making use of the commutativity $BC = CB$, we rearrange the series 
\[
\sum\limits_{k=0} ^{\infty} (t (B + C)) ^{k} x / k!
\]
so that we obtain $\left(\sum\limits_{k=0}^{\infty}(tB)^{k}/k!\right)\left(\sum\limits_{k=0}^{\infty}(tC)^{k}x/k!\right)$ as in the case of numerical series $\sum\limits_{k=0}^{\infty}(t(b+c))^{k}/k!$. 
\end{proof}

\begin{corollary}\label{cor:IX.6.Corollary.3} For every $x \in X$, 
\begin{equation}\label{eq:IX.6.4}
\lim\limits _{h \downarrow 0} h ^{-1} (\exp (h B) - I) x = B x,
\end{equation}
and so, by making use of the semi-group property 
\begin{equation}\label{eq:IX.6.5}
\exp ((t + h) B) = \exp (t B) \exp (h B)
\end{equation}
proved above, we obtain 
\begin{equation}\label{eq:IX.6.6}
D _{t} \exp (t B) x = \exp (t B) B x = B \exp (t B) x.
\end{equation}
\end{corollary}

\begin{proof}
 For any continuous semi-norm $\phi$ on $X$, we obtain, as above 
\[
\begin{aligned}
p (h ^{-1} (\exp (h B) - I) x - B x)& \leqq \sum\limits_{k=2} ^{\infty} h ^{k-1} p (B ^{k} x) / k! \leqq q (x) \sum\limits_{k=2} ^{\infty} h ^{k-1} / k!
\\
& \text{ which surely tends to } 0 \text{ as } h \downarrow 0. \end{aligned}
\]
\end{proof}


\yossection[7. The Representation and the Characterization of Equi-continuous Semi-groups]{7}{The Representation and the Characterization of Equi-continuous Semi-groups of Class \texorpdfstring{$(C_{0})$}{(C0)} in Terms of the Corresponding Infinitesimal Generators}
\label{sec:IX.7}

We shall prove the following fundamental 

\begin{theoremu} Let $X$ be a locally convex, sequentially complete, linear topological space. Suppose $A$ is a linear operator with domain $D(A)$ dense in $X$ and range $R(A)$ in $X$ such that the resolvent $R(n; A) = (nI - A)^{-1} \in L(X, X)$ exists for $n = 1, 2, \ldots$. Then $A$ is the infinitesimal generator of a uniquely determined equi-continuous semi-group of class $(C_0)$ iff the operators $\{(I - n^{-1}A)^{-m}\}$ are equi-continuous in $n = 1, 2, \ldots$ and $m = 0, 1, \ldots$. 
\end{theoremu}

\begin{proof}
 The ``only if'' part is already proved. We shall prove the ``if'' part. 

Setting 
\begin{equation}\label{eq:IX.7.1}
J _{n} = (I - n ^{-1} A) ^{-1},
\end{equation}
we shall prove 
\begin{equation}\label{eq:IX.7.2}
\lim\limits _{n \rightarrow \infty} J _{n} x = x \quad \text{ for every } \quad x \in X.
\end{equation}

In fact, we have, for $x \in D(A)$, $A J_n x = J_n A x = n(J_n - I)x$ and so $J_n x - x = n^{-1} J_n(Ax)$ tends to $0$ as $n \to \infty$, in virtue of the equi-boundedness of $\{J_n(Ax)\}$ in $n = 1, 2, \ldots$. Since $D(A)$ is dense in $X$ and $\{J_n\}$ is equi-continuous in $n$, it follows that $\lim\limits_{n \to \infty} J_n x = x$ for every $x \in X$. 


Put 
\begin{equation}\label{eq:IX.7.3}
T _{t} ^{(n)} = \exp (t A J _{n}) = \exp (t n (J _{n} - I)) = \exp (- n t) \exp (n t J _{n}), t \geqq 0.
\end{equation}

Since $\{J_n^k\}$ is equi-bounded in $n$ and $k$, the $\exp(t n J_n)$ can be defined, and we have, as in \eqref{eq:IX.7.2} of the preceding section, 
\[
p \left(\exp (n t J _{n}) x\right) \leqq \sum\limits_{k=0} ^{\infty} (n t) ^{k} (k!) ^{-1} p \left(J _{n} ^{k} \cdot x\right) \leqq \exp (n t) \cdot q (x).
\]

Consequently, the operators $\{T_t^{(n)}\}$ are equi-continuous in $t \geqq 0$ and $n = 1, 2, \ldots$ in such a way that 
\begin{equation}\label{eq:IX.7.4}
p \left(T _{t} ^{(n)} x\right) \leqq q (x).
\end{equation}

We next remark that $J_{n}J_{m}=J_{m}J_{n}$ for $n, m>0$. Thus $J_{n}$ is commutative with $T_{t}^{(m)}$. Hence, by $D_{t}T_{t}^{(n)}x=AJ_{n}T_{t}^{(n)}x=T_{t}^{(n)}AJ_{n}x$, proved in the preceding section, 
\begin{equation}\label{eq:IX.7.5}
\begin{aligned}
&\phi (T _{t} ^{(n)} x - T _{t} ^{(m)} x)  = \phi \left(\int_{0} ^{t} D _{s} (T _{t-s} ^{(m)} T _{s} ^{(n)} x) \dd s\right) = \phi \left(\int_{0} ^{t} T _{t-s} ^{(m)} T _{s} ^{(n)} (A J _{n} - A J _{m}) x \dd s\right).
\end{aligned}
\end{equation}

Hence, if $x \in D(A)$, there exists a continuous semi-norm $\tilde{q}$ on $X$ such that 
\[
p \left(T _{t} ^{(n)} x - T _{t} ^{(m)} x\right) \leqq \int_{0} ^{t} \tilde {q} \left(\left(A J _{n} - A J _{m}\right) x\right) \dd s = t \tilde {q} \left(\left(J _{n} A - J _{m} A\right) x\right).
\]

Therefore, by \eqref{eq:IX.7.2}, we have proved that $\lim\limits_{n,m\to \infty}p(T_t^{(n)}x - T_t^{(m)}x) = 0$ uniformly in $t$ when $t$ varies on every compact interval. Since $D(A)$ is dense in the sequentially complete space $X$, and since the operators $\{T^{(n)}\}$ are equi-continuous in $t\geqq 0$ and in $n$, we see that $\lim\limits_{n\to \infty}T_t^{(n)}x = T_tx$ exists for every $x\in X$ and $t\geqq 0$ uniformly in $t$ on every compact interval of $t$. Thus the operators $\{T_t\}$ are equi-continuous in $t\geqq 0$, and from the uniform convergence in $t$, $T_tx$ is continuous in $t\geqq 0$. 

We next prove that $T_{t}$ satisfies the semi-group property $T_{t}T_{s} = T_{t+s}$. Since $T_{t+s}^{(n)} = T_{t}^{(n)}T_{s}^{(n)}$, we have 
\[
\begin{aligned}
&\phi (T _{t+s} x - T _{t} T _{s} x)\leqq \phi (T _{t+s} x - T _{t+s} ^{(n)} x) + \phi (T _{t+s} ^{(n)} x - T _{t} ^{(n)} T _{s} ^{(n)} x) \\
&\quad + \phi (T _{t} ^{(n)} T _{s} ^{(n)} x - T _{t} ^{(n)} T _{s} x) + \phi (T _{t} ^{(n)} T _{s} x - T _{t} T _{s} x) \leqq \phi (T _{t+s} x - T _{t+s} ^{(n)} x) + q (T _{s} ^{(n)} x - T _{s} x) \\
&\quad + \phi ((T _{t} ^{(n)} - T _{t}) T _{s} x) \rightarrow 0 \text{ as } n \rightarrow \infty .
\end{aligned}
\]

Thus $p(T_{t+s}x - T_tT_s x) = 0$ for every continuous semi-norm $p$ on $X$, which proves that $T_{t+s} = T_tT_s$. 

Let $\hat{A}$ be the infinitesimal generator of this equi-continuous semigroup $\{T_t\}$ of class $(C_0)$. We have to show that $\hat{A} = A$. Let $x \in D(A)$. Then $\lim\limits_{n \to \infty} T_t^{(n)} A J_n x = T_t A x$ uniformly in $t$ on every compact interval of $t$. For, we have, by \eqref{eq:IX.7.4}, 
\[
\begin{aligned} p (T _{t} A x - T _{t} ^{(n)} A J _{n} x) & \leqq p (T _{t} A x - T _{t} ^{(n)} A x) + p (T _{t} ^{(n)} A x - T _{t} ^{(n)} A J _{n} x) \\ & \leqq p ((T _{t} - T _{t} ^{(n)}) A x) + q (A x - J _{n} A x) \end{aligned}
\]
which tends to $0$ as $n \to \infty$, since $\lim\limits_{n \to \infty} J_n A x = A x$. Hence 
\[
\begin{aligned}
&T _{t} x - x= \lim\limits _{n \rightarrow \infty} (T _{t} ^{(n)} x - x) = \lim\limits _{n \rightarrow \infty} \int_{0} ^{t} T _{s} ^{(n)} A J _{n} x \dd s = \int_{0} ^{t} (\lim\limits _{n \rightarrow \infty} T _{s} ^{(n)} A J _{n} x) \dd s = \int_{0} ^{t} T _{s} A x \dd s
\end{aligned}
\]
so that $\lim\limits_{t\downarrow 0}t^{-1}(T_t x - x) = \lim\limits_{t\downarrow 0}t^{-1}\int_0^t T_sAxds$ exists and equals $A x$. We 
have thus proved that $x \in D(A)$ implies $x \in D(\hat{A})$ and $Ax = \hat{A}x$, i.e. $\hat{A}$ is an extension of $A$. $\hat{A}$ being the infinitesimal generator of the semi-group $T_{t}$, we know that, for $n > 0$, $(nI - \hat{A})$ maps $D(\hat{A})$ onto $X$ in one-one way. But, by hypothesis, $(nI - A)$ also maps $D(A)$ onto $X$ in one-one way. Therefore the extension $\hat{A}$ of $A$ must coincide with $A$. 

Finally, the uniqueness of the semi-group $T_{t}$ is proved as follows. Let $\tilde{T}_{t}$ be an equi-continuous semi-group of class $(C_0)$ whose infinitesimal generator is precisely $A$. We construct the semi-group $T_{t}^{(n)}$. Since $A$ is commutative with $\tilde{T}_{t}$, we see that $AJ_{n}$ and $T_{s}^{(n)}$ are commutative with $\tilde{T}_{t}$. Then, for $x \in D(A)$, we obtain, as in \eqref{eq:IX.7.5}, 
\begin{equation}\label{eq:IX.7.6}
\begin{aligned}
&p (T _{t} ^{(n)} x - \tilde {T} _{t} x)  = p \left(\int_{0} ^{t} D _{s} (\tilde {T} _{t-s} T _{s} ^{(n)} x) \dd s\right) = p \left(\int_{0} ^{t} - \tilde {T} _{t-s} T _{s} ^{(n)} (A - A J _{n}) x \dd s\right).
\end{aligned}
\end{equation}

Thus, by virtue of $\lim\limits_{n\to \infty}AJ_nx = Ax$ for all $x\in D(A)$, we prove $\lim\limits_{n\to \infty}T_t^{(n)}x = \tilde{T}_t x$ for all $x\in X$ similarly to the above proof of the existence of $\lim\limits_{n\to \infty}T_t^{(n)}x,x\in X$. 

Therefore, $T_{t}x = \tilde{T}_{t}x$ for all $x \in X$, that is, $T_{t} = \tilde{T}_{t}$.
\end{proof}

\begin{remarku} The above proof shows that, if $A$ is the infinitesimal generator of an equi-continuous semi-group $T_{t}$ of class $(C_{0})$, then 
\begin{equation}\label{eq:IX.7.7}
T _{t} x = \lim\limits _{n \rightarrow \infty} \exp (t A (I - n ^{-1} A) ^{-1}) x, x \in X,
\end{equation}
\end{remarku}

and the convergence in \eqref{eq:IX.7.7} is uniform in $t$ on every compact interval of $t$. This is the representation theorem for semi-groups. 

\begin{corollary}\label{cor:IX.7.Corollary.1} If $X$ is a $B$-space, then the condition of the Theorem reads: $D(A)^{a} = X$ and the resolvent $(I - n^{-1}A)^{-1}$ exists such that 
\begin{equation}\label{eq:IX.7.8}
\left\| (I - n ^{-1} A) ^{-m} \right\| \leqq C \quad (n = 1, 2, \dots ; m = 1, 2, \dots)
\end{equation}
with a positive constant $C$ which is independent of $n$ and $m$. In particular, for the case of a contraction semi-group, the condition reads: $D(A)^{a} = X$ and the resolvent $(I - n^{-1}A)^{-1}$ exists such that 
\begin{equation}\label{eq:IX.7.9}
\left\| (I - n ^{-1} A) ^{-1} \right\| \leqq 1 (n = 1, 2, \dots).
\end{equation}
\end{corollary}

\begin{remarku} The above result \eqref{eq:IX.7.9} was obtained independently by E. \textsc{Hille} \cite{hille98} and by K. \textsc{Yosida} \cite{yosida348}. The result was extended by W. \textsc{Feller} \cite{erdelyi56}, R. S. \textsc{Phillips} \cite{phillips253} and I. \textsc{Miyadera} \cite{miyadera204} and the extension is given in the form \eqref{eq:IX.7.8}. It is to be noted that in condition \eqref{eq:IX.7.8} and \eqref{eq:IX.7.9}, we may replace $(n=1,2,\ldots)$ by (for all sufficiently large $n$). The extension of the semi-group theory to locally convex linear topological spaces, as given in the present book, is suggested by L. \textsc{Schwartz} \cite{schwartz296}. 
\end{remarku}

\begin{corollary}\label{cor:IX.7.Corollary.2} Let $X$ be a $B$-space, and $\{T_{t}; t \geqq 0\}$ be a family of bounded linear operators $L(X, X)$ such that 
\begin{align}
T _{t} T _{s}& = T _{t+s} (t, s \geqq 0), T _{0} = I, \label{eq:IX.7.10} \\
\slim _{t \downarrow 0} T _{t} x& = x \text{ for all } x \in X, \label{eq:IX.7.11} \\
\left\| T _{t} \right\|& \leqq M e ^{\beta t} \text{ for all } t \geqq 0, \text{ where } M > 0 \text{ and } \beta \geqq 0 \text{ are } \label{eq:IX.7.12}
\end{align}
independent of $t$. 
\end{corollary}


Then $(A - \beta I)$ is the infinitesimal generator of the equi-continuous semi-group $S_{t} = e^{-\beta t}T_{t}$ of class $(C_0)$, where $A$ is the operator defined through $Ax = s\lim\limits_{t\downarrow 0}t^{-1}(T_t - I)x$. Thus, by Corollary~\ref{cor:IX.7.Corollary.1}, we see that a closed linear operator $A$ with $D(A)^{a} = X$ and $R(A)\subseteq X$ is the infinitesimal generator of a semi-group $T_{t}$ satisfying \eqref{eq:IX.7.10}, \eqref{eq:IX.7.11} and \eqref{eq:IX.7.12} iff the resolvent $(I - n^{-1}(A - \beta I))^{-1}$ exists such that 
$\left\|(I-n^{-1}(A-\beta I))^{-m}\right\|\leqq M$ (for $m=1,2,\ldots$ and all large $n$). This condition may be rewritten as 

\begin{equation}\tag{13'}\label{eq:IX.7.13'}
\left\| (I - n ^{-1} A) ^{-m} \right\| \leqq M (1 - n ^{-1} \beta) ^{-m} (\text{ for } m = 1, 2, \dots \text{ and all large } n).
\end{equation}

In particular, for those semi-groups $T_{t}$ satisfying \eqref{eq:IX.7.10}, \eqref{eq:IX.7.11} and 
\begin{equation}\label{eq:IX.7.14}
\left\| T _{t} \right\| \leqq e ^{\beta t} \quad \text{ for all } \quad t \geqq 0,
\end{equation}
condition \eqref{eq:IX.7.13'} may be replaced by 
\begin{equation}\tag{13''}\label{eq:IX.7.13}
\left\| (I - n ^{-1} A) ^{-1} \right\| \leqq (1 - n ^{-1} \beta) ^{-1} \quad (\text{ for all large } n)
\end{equation}

An application of the representation theorem to the proof of Weierstrass' polynomial approximation theorem. Consider the semi-group $T_{t}$ defined by $(T_{t}x)(s) = x(t + s)$ on $C[0,\infty]$. The representation theorem gives 
\[
(T _{t} x) (s) = x (t + s) = \slim\limits_{n \rightarrow \infty} \exp (t A J _{n} x) (s) = \slim\limits_{n \rightarrow \infty} \sum\limits_{m=0} ^{\infty} \frac {t ^{m}}{m !} (A J _{n}) ^{m} x (s),
\]
and the above \slim is uniform in $t$ on any compact interval of $t$. From this result we can derive Weierstrass' polynomial approximation theorem. Let $z(s)$ be a continuous function on the closed interval [$0$, $1$]. Let $x(s) \in C[0, \infty]$ be such that $x(s) = z(s)$ for $s \in [0, 1]$. Put $s = 0$ in the above representation of $x(t + s)$. Then we obtain 
\[
(T _{t} x) (0) = x (t) = \slim _{n \to \infty} \sum\limits_{m=0} ^{\infty} t ^{m} (A J _{n}) ^{m} x (0) / m! \quad \text{ in } \quad C [ 0, 1 ],
\]
uniformly in $t$ on $[0, 1]$. Hence $z(t)$ is a uniform limit on $[0, 1]$ of a sequence of polynomials in $t$. 


\yossection{8}{Contraction Semi-groups and Dissipative Operators}
\label{sec:IX.8}

G. \textsc{Lumer} and R. S. \textsc{Phillips} have discussed contraction semi-groups by virtue of the notion of semi-scalar product. The infinitesimal generator of such a semi-group is dissipative in their terminology. 

\begin{propositionu} (\textsc{Lumer}). To each pair $\{x, y\}$ of a complex (or real) normed linear space $X$, we can associate a complex (or real) number $[x, y]$ such that 
\begin{equation}\label{eq:IX.8.1}
\begin{gathered}
{}[ x + y, z ] = [ x, z ] + [ y, z ], \quad [ \lambda x, y ] = \lambda [ x, y ], \quad [ x, x ] = \| x \| ^{2},
\\
| [ x, y ] | \leqq \| x \| \cdot \| y \|.
\end{gathered}
\end{equation}
$[x, y]$ is called a semi-scalar product of the vectors $x$ and $y$. 
\end{propositionu}

\begin{proof}
 According to Corollary~\hyperref[cor:IV.6.Corollary.2]{$2$} of Theorem~\ref{thm:IV.6.Theorem.1} in Chapter~\ref{chap:IV}, \ref{sec:IV.6}, there exists, for each $x_0 \in X$, at least one (and let us choose exactly one) bounded linear functional $f_{x_0} \in X'$ such that $\|f_{x_0}\| = \|x_0\|$ and $\langle x_0, f_{x_0} \rangle = \|x_0\|^2$. Then clearly 
\begin{equation}\label{eq:IX.8.2}
[ x, y ] = \langle x, f _{y} \rangle
\end{equation}
defines a semi-scalar product. 
\end{proof}

\begin{definitionu} Let a complex (or real) $B$-space $X$ be endowed with a semi-scalar product $[x, y]$. A linear operator $A$ with domain $D(A)$ and range $R(A)$ both in $X$ is called dissipative (with respect to $[x, y]$) if 
\begin{equation}\label{eq:IX.8.3}
\operatorname{Re} [ A x, x ] \leqq 0 \quad \text{ whenever } \quad x \in D (A).
\end{equation}
\end{definitionu}

\begin{examplen} Let $X$ be a Hilbert space. Then a symmetric operator $A$ such that $(Ax, x) \leqq 0$ is surely dissipative with respect to the semi-scalar product $[x, y] = (x, y)$, where $(x, y)$ is the ordinary scalar product.

\end{examplen}
\begin{theoremn}[\textsc{Phillips} and \textsc{Lumer}] Let $A$ be a linear operator with domain $D(A)$ and range $R(A)$ both in a complex (or real) $B$-space $X$ such that $D(A)^{a} = X$. Then $A$ generates a contraction semi-group of class $(C_{0})$ in $X$ iff $A$ is dissipative (with respect to any semi-scalar product $[x, y]$) and $R(I - A) = X$.

\end{theoremn}
\begin{proof}
 The ``if'' part. Let $A$ be dissipative and $\lambda > 0$. Then the inverse $(\lambda I - A)^{-1}$ exists and $\|(\lambda I - A)^{-1}y\| \leqq \lambda^{-1}\|y\|$ when $y \in D((\lambda I - A)^{-1})$. 

For, if $y = \lambda x - Ax$, then
\begin{equation}\label{eq:IX.8.4}
\lambda\|x\|^2=\lambda[x,x]\leqq\operatorname{Re}(\lambda[x,x]-[Ax,x])=\operatorname{Re}[y,x]\leqq\|y\|\cdot\|x\|,
\end{equation}

since $A$ is dissipative. By hypothesis, $R(I - A) = X$ so that $\lambda = 1$ is in the resolvent set $\varrho (A)$ of $A$, and we have $\|R(1;A)\|\leqq 1$ by \eqref{eq:IX.8.4}. If $|\lambda -1| < 1$, then the resolvent $R(\lambda ;A)$ exists and is given by 
\[
\begin{aligned}
&R (\lambda ; A)  = R (1; A) \left(I + (\lambda - 1) R (1; A)\right) ^{-1} = R (1; A) \cdot \sum\limits_{n=0} ^{\infty} \big ((1 - \lambda) R (1; A) \big) ^{n},
\end{aligned}
\]

(see Theorem~\ref{thm:VIII.2.Theorem.1} in Chapter~\ref{chap:VIII}, \ref{sec:VIII.2}). Moreover, \eqref{eq:IX.8.4} implies that $\|R(\lambda; A)\| \leqq \lambda^{-1}$ for $\lambda > 0$ with $|\lambda - 1| < 1$. Hence, again by 
\[
R (\mu ; A) = R (\lambda ; A) (I + (\mu - \lambda) R (\lambda ; A)) ^{-1},
\]
which is valid for $\mu > 0$ with $|\mu - \lambda| \|R(\lambda; A)\| < 1$, we prove the existence of $R(\mu; A)$ and $\|R(\mu; A)\| \leqq \mu^{-1}$. Repeating the process, we see that $R(\lambda; A)$ exists for all $\lambda > 0$ and satisfies the estimate $\|R(\lambda; A)\| \leqq \lambda^{-1}$. As $D(A)$ is dense by hypothesis, it now follows, from Corollary~\hyperref[cor:IX.7.Corollary.1]{$1$} of the preceding section, that $A$ generates a contraction semi-group of class $(C_0)$. 


The ``only if'' part. Suppose $\{T_{t}; t \geqq 0\}$ is a contraction semi-group of class $(C_{0})$. Then 
\begin{gather}\operatorname{Re} \left[ T _{t} x - x, x \right] = \operatorname{Re} \left[ T _{t} x, x \right] - \| x \| ^{2} \leqq \| T _{t} x \| \cdot \| x \| - \| x \| ^{2} \leqq 0. x \in D (A) T, \notag \\\operatorname{Re} [ A x, x ] = \lim\limits _{t \downarrow 0} \operatorname{Re} \left\{t ^{-1} \left[ T _{t} x - x, x \right] \right\} \leqq 0 \notag
\end{gather}

Moreover, we know that $R(I - A) = D(R(1; A)) = X$, since $A$ is the infinitesimal generator of a contraction semi-group of class $(C_0)$.
\end{proof}

\begin{corollaryu} If $A$ is a densely defined closed linear operator such that $D(A)$ and $R(A)$ are both in a $B$-space $X$ and if $A$ and its dual operator $A'$ are both dissipative, then $A$ generates a contraction semi-group of class $(C_{0})$. 
\end{corollaryu}

\begin{proof}
 It suffices to show that $R(I - A) = X$. But, since $(I - A)^{-1}$ is closed with $A$ and continuous, $R(I - A)$ is a closed linear subspace of $X$. Thus $R(I - A) \neq X$ implies the existence of a non-trivial $x' \in X'$ such that $\langle (x - Ax), x' \rangle = 0$ for all $x \in D(A)$. 

Hence $x' - A'x' = 0$, contrary to the dissipativity of $A'$ and $x' \neq 0$. 
\end{proof}

\begin{remarku} For further details concerning dissipative operators, see G. \textsc{Lumer}-R. S. \textsc{Phillips} \cite{lumer183}. See also T. \textsc{Kato} \cite{kato139}. 
\end{remarku}


\yossection{9}{Equi-continuous Groups of Class \texorpdfstring{$(C_0)$}{(C0)}. Stone's Theorem}
\label{sec:IX.9}

\begin{definitionu} An equi-continuous semi-group $\{T_t\}$ of class $(C_0)$ is called an equi-continuous group of class $(C_0)$, if there exists an equi-continuous semi-group $\{\hat{T}_t\}$ of class $(C_0)$ with the condition: 

If we define $S_{t}$ by $S_{t} = T_{t}$ for $t \geqq 0$ and $S_{-t} = \hat{T}_{t}$ for $t \geqq 0$, then the family of operators $S_{t}, -\infty < t < \infty$, has the group property 
\end{definitionu}
\begin{equation}\label{eq:IX.9.1}
S _{t} S _{s} = S _{t+s} (- \infty < t, s < \infty), S _{0} = I.
\end{equation}

\begin{theoremu} Let $X$ be a locally convex, sequentially complete linear topological space. Suppose $A$ is a linear operator with domain $D(A)$ dense in $X$ and range $R(A)$ in $X$. Then $A$ is the infinitesimal generator of an equi-continuous group of class $(C_{0})$ of operators $S_{t} \in L(X, X)$ iff the operators $(I - n^{-1}A)^{-m}$ are everywhere defined and equi-continuous in $n = \pm 1, \pm 2, \ldots$ and in $m = 1, 2, \ldots$. 
\end{theoremu}

\begin{proof}
 The ``only if'' part. Let $T_{t}=S_{t}$ for $t\geqq0$ and $\hat{T}_{t}=S_{-t}$ for $t\geqq0$. Let $A$ and $\hat{A}$ be the infinitesimal generator of $\hat{T}_{t}$ and $T_{t}$, respectively. We have to show that $\hat{A}=-A$. If $x\in D(\hat{A})$, then, by putting $x_{h}=h^{-1}(\hat{T}_{h}-I)x$ and making use of the equi-continuity of $T_{h}$, 
\[
\begin{aligned}
&p (T _{h} x _{h} - \hat {A} x)  \leqq p (T _{h} x _{h} - T _{h} \hat {A} x) + p ((T _{h} - I) A x) \leqq q (x _{h} - \hat {A} x) + p ((T _{h} - I) \hat {A} x)
\end{aligned}
\]
where $p$ and $q$ are continuous semi-norms on $X$ such that, for a given $p$, we can choose a $q$ satisfying the above inequality for all $h \geqq 0$ and all $x \in D(\hat{A})$ simultaneously. Thus $\lim\limits_{h \downarrow 0} T_h x_h = \hat{A} x$, and so $x \in D(\hat{A})$ implies $\hat{A} x = \lim\limits_{h \downarrow 0} T_h (h^{-1}(\hat{T}_h - I)) x = \lim\limits_{h \downarrow 0} h^{-1}(I - T_h) x = -Ax$. Hence $-A$ is an extension of $\hat{A}$. In the same way, we can prove that $\hat{A}$ is an extension of $-A$. Therefore $\hat{A} = -A$. 

The ``if'' part. Define, for $t \geqq 0$, 
\begin{align}
T _{t} x& = \lim\limits _{n \rightarrow \infty} T _{t} ^{(n)} x = \lim\limits _{n \rightarrow \infty} \exp (t A (I - n ^{-1} A) ^{-1}) x, \notag \\
\hat {T} _{t} x& = \lim\limits _{n \rightarrow \infty} \hat {T} _{t} ^{(n)} x = \lim\limits _{n \rightarrow \infty} \exp (t \hat {A} (I - n ^{-1} \hat {A}) ^{-1}) x, \text{ where } \hat {A} = - A. \notag
\end{align}

Then $T_{t}$ and $\hat{T}_{t}$ are both equi-continuous semi-groups of class $(C_{0})$. We have 
\[
\begin{aligned} p (T _{t} \hat {T} _{t} x - T _{t} ^{(n)} \hat {T} _{t} ^{(n)} x) & \leqq p (T _{t} \hat {T} _{t} x - T _{t} ^{(n)} \hat {T} _{t} x) + p (T _{t} ^{(n)} \hat {T} _{t} x - T _{t} ^{(n)} \hat {T} _{t} ^{(n)} x) \\ & \leqq p ((T _{t} - T _{t} ^{(n)}) \hat {T} _{t} x) + q (\hat {T} _{t} x - \hat {T} _{t} ^{(n)} x) \end{aligned}
\]
by the equi-continuity of $\{T_t^{(n)}\}$ in $n$ and $t \geqq 0$. Thus we have $\lim\limits_{n \to \infty} T_t^{(n)} \hat{T}_t^{(n)} x = T_t \hat{T}_t x$. We have incidentally proved the equi-continuity of $T_t \hat{T}_t$ by the equi-continuity of $T_t^{(n)} \hat{T}_t^{(n)}$. On the other hand, we have $T_t^{(n)} \hat{T}_s^{(m)} = \hat{T}_s^{(m)} T_t^{(n)}$ by the commutativity of $(I - n^{-1}A)^{-1}$ and $(I - m^{-1}A)^{-1}$. Thus $(T_t^{(n)} \hat{T}_t^{(n)}) (T_s^{(n)} \hat{T}_s^{(n)}) = T_{t+s}^{(n)} \hat{T}_{t+s}^{(n)}$, and so $(T_{t}\hat{T}_{t})$ with $t \geqq 0$ enjoys the semi-group property $(T_{t}\hat{T}_{t})(T_{s}\hat{T}_{s}) = T_{t+s}\hat{T}_{t+s}$. Hence $\{T_{t}\hat{T}_{t}\}$ is an equi-continuous semi-group of class $(C_0)$. If $x \in D(\hat{A}) = D(A)$, then 
\[
\begin{aligned}
&\lim\limits _{h \downarrow 0} h ^{-1} (T _{h} \hat {T} _{h} - I) x  = \lim\limits _{h \downarrow 0} T _{h} h ^{-1} (\hat {T} _{h} - I) x + \lim\limits _{h \downarrow 0} h ^{-1} (T _{h} - I) x = \hat {A} x + A x = 0,
\end{aligned}
\]
so that the infinitesimal generator $A_{1}$ of $\{T_t\hat{T}_t\}$ is $0$ at every $x \in D(\hat{A})$. Since $(I - A_1)$ is the inverse of a continuous linear operator $(I - A_1)^{-1} \in L(X, X)$, we see that $A_{1}$ must be closed. Hence $A_{1}$ must be = $0$, in virtue of the fact that $A_{1}$ vanishes on a dense subset $D(\hat{A}) = D(A)$ of $X$. Therefore $(T_t\hat{T}_t)x = \lim\limits_{n\to \infty}\exp (t\cdot 0\cdot (I - n^{-1}\cdot 0)^{-1})x = x$, that is, $T_{t}\hat{T}_{t} = I$. We have thus proved the group property of $S_{t}, -\infty < t < \infty$, where $S_{t} = T_{t}$ and $S_{-t} = \hat{T}_t$ for $t\geqq 0$. 
\end{proof}

\begin{corollary}\label{cor:IX.9.Corollary.1} For the case of a $B$-space $X$, the condition of the Theorem reads: $D(A)^{a} = X$ and the resolvent $(I - n^{-1}A)^{-1}$ exists such that 
\begin{equation}\label{eq:IX.9.2}
\left\|(I - n ^{-1} A)^{-m}\right\| \leqq M \text{ (for } m = 1, 2, \ldots \text{ and all large } |n|, n \geqq 0).
\end{equation}
For the group $S_{t}$ satisfying $\left\|S_{t}\right\| \leqq M e^{\beta |t|} (\beta \geqq 0)$ for all $t, -\infty < t < \infty$, the condition reads: $D(A)^{a} = X$ and the resolvent $(I - n^{-1}A)^{-1}$ exists such that 
\begin{equation}\label{eq:IX.9.3}
\left\| (I - n ^{-1} A) ^{-m} \right\| \leqq M (1 - | n ^{-1} | \beta) ^{-m} \text{ (for } m = 1, 2, \ldots \text{ and all large } |n|, n \leqq 0).
\end{equation}
\end{corollary}


For the particular case $\| S_t \| \leqq e^{\beta |t|}$ for all $t, -\infty < t < \infty$, the condition reads: $D(A)^a = X$ and the resolvent $(I - n^{-1}A)^{-1}$ exists such that 
\begin{equation}\label{eq:IX.9.4}
\left\| (I - n ^{-1} A) ^{-1} \right\| \leqq (1 - | n ^{-1} | \beta) ^{-1} (\text{ for all large } | n |, n \geqq 0).
\end{equation}

\begin{proof}
 As the proof of Corollary~\ref{cor:IX.9.Corollary.1} and Corollary~\ref{cor:IX.7.Corollary.2} in Chapter~\ref{chap:IX}, \ref{sec:IX.7}. \end{proof}

\begin{corollary}[Stone's theorem]\label{cor:IX.9.Corollary.2} Let $U_{t}, -\infty < t < \infty$, be a group of unitary operators of class $(C_{0})$ in a Hilbert space $X$. Then the infinitesimal generator $A$ of $U_{t}$ is $\sqrt{-1}$ times a self-adjoint operator $H$. 
\end{corollary}

\begin{proof}
 We have $(U_{t}x,y) = (x,U_{t}^{-1}y) = (x,U_{-t}y)$ and so, by differentiation, 
\[
(A x, y) = (x, - A y) \text{ whenever } x \text{ and } y \in D (A) \]

Thus $-iA = H$ is symmetric. $A$ being the infinitesimal generator of $U_{t}$, $(I - n^{-1}A)^{-1} = (I - n^{-1}iH)^{-1}$ must be a bounded linear operator such that $\left\|(I - n^{-1}iH)^{-1}\right\| \leqq 1$ for $n = \pm 1, \pm 2, \ldots$. Hence, taking the case $n = \pm 1$, we see that the Cayley transform of $H$ is unitary. This proves that $H$ is self-adjoint. 
\end{proof}

\begin{remarku} If $A$ is of the form $A = \sqrt{-1} \cdot H$, where $H$ is a self-adjoint operator in a Hilbert space $X$, then condition \eqref{eq:IX.9.4} of Corollary~\ref{cor:IX.9.Corollary.1} is surely satisfied, as may be proved by the theory of Cayley's transform. Therefore, $A$ is the infinitesimal generator of a group of contraction operators $U_t$ in $X$. It is easy to see that such $U_t$ is unitary. For, a contraction operator $U_t$ on a Hilbert space $X$ into $X$ must be unitary, if $U_t^{-1} = U_{-t}$ is also a contraction operator on $X$ into $X$. 
\end{remarku}


\yossection{10}{Holomorphic Semi-groups}
\label{sec:IX.10}

We shall introduce an important class of semi-groups, namely, semi-groups $T_{t}$ which can, as functions of the parameter $t$, be continued holomorphically into a sector of the complex plane containing the positive $t$-axis. We first prove 

\begin{lemmau} Let $X$ be a locally convex, sequentially complete, linear topological space. Let $\{T_{t}; t \geqq 0\} \subseteq L(X, X)$ be an equi-continuous semigroup of class $(C_{0})$. Suppose that, for all $t > 0$, $T_{t}X \subseteq D(A)$, the domain of the infinitesimal generator $A$ of $T_{t}$. Then, for any $x \in X$, $T_{t}x$ is infinitely differentiable in $t > 0$ and we have 
\begin{equation}\label{eq:IX.10.1}
T _{t} ^{(n)} x = (T _{t/n} ^{\prime}) ^{n} x \quad \text{ for all } \quad t > 0,
\end{equation}
where $T_{t}^{\prime} = D_{t}T_{t}, T_{t}^{\prime \prime} = D_{t}T_{t}^{\prime}, \ldots, T_{t}^{(n)} = D_{t}T_{t}^{(n-1)}$. 
\end{lemmau}

\begin{proof}
 If $t > t_0 > 0$, then $T_t'x = AT_tx = T_{t-t_0}AT_{t_0}x$ by the commutativity of $A$ and $T_s, s \geqq 0$. Thus $T_t'X \subseteq T_{t-t_0}X \subseteq D(A)$ when $t > 0$, and so $T_t''x$ exists for all $t > 0$ and $x \in X$. Since $A$ is a closed linear operator, we have 
\[ \begin{aligned}
&T _{t} ^{\prime \prime} x= D _{t} (A T _{t}) x = A \cdot \lim\limits _{n \rightarrow \infty} \cdot n (T _{t+1/n} - T _{t}) x = A (A T _{t}) x = A T _{t/2} A T _{t/2} x = (T _{t/2} ^{\prime}) ^{2} x.
\end{aligned}
\]

Repeating the same argument, we obtain \eqref{eq:IX.10.1}. 
\end{proof}

Let $X$ be a locally convex, sequentially complete, complex linear topological space. Let $\{T_{t}; t \geqq 0\} \subseteq L(X, X)$ be an equi-continuous semi-group of class $(C_{0})$. For such a semi-group, we consider the following three conditions: 

(I) For all $t > 0$, $T_{t}x \in D(A)$, and there exists a positive constant $C$ such that the family of operators $\{(CtT_t)^n\}$ is equi-continuous in $n \geqq 0$ and $t, 0 < t \leqq 1$. 

(II) $T_{t}$ admits a weakly holomorphic extension $T_{\lambda}$ given locally by 
\begin{equation}\label{eq:IX.10.2}
T _{\lambda} x = \sum\limits_{n=0} ^{\infty} (\lambda - t) ^{n} T _{t} ^{(n)} x / n! \quad \text{ for } \quad | \arg \lambda | < \tan^{-1} (C e ^{-1}),
\end{equation}
the family of operators $\{e^{-\lambda}T_{\lambda}\}$ is equi-continuous in $\lambda$ for 
\begin{equation}\label{eq:IX.10.3}
\left| \arg \lambda \right| < \tan^{-1} \left(2 ^{-k} C e ^{-1}\right) \text{ with some positive constant } k.
\end{equation}

(III) Let $A$ be the infinitesimal generator of $T_{t}$. Then there exists a positive constant $C_{1}$ such that the family of operators $\{(C_{1}\lambda R(\lambda ;A))^{n}\}$ is equi-continuous in $n\geqq 0$ and in $\lambda$ with $\operatorname{Re}(\lambda)\geqq 1 + \varepsilon ,\varepsilon >0$. 

We can prove 

\begin{theoremu} The three conditions (I), (II) and (III) are mutually equivalent. 
\end{theoremu}

\begin{proof}
 The implication (I) $\rightarrow$ (II). Let $p$ be any continuous semi-norm on $X$. Then, by hypothesis, there exists a continuous semi-norm $q$ on $X$ such that $p((tT_t')^n x) \leqq C^{-n} q(x)$, $0 < C$ for $1 \geqq t > 0$, $n \geqq 0$ and $x \in X$. Hence, by \eqref{eq:IX.10.1}, we obtain, for any $t > 0$, 
\[ \begin{aligned} \phi ((\lambda - t) ^{n} T _{t} ^{(n)} x / n!) & \leqq \frac {| \lambda - t | ^{n}}{t ^{n}} \frac {n ^{n}}{n !} \frac {1}{C ^{n}} \phi \left(\left(\frac {t}{n} C T _{t/n} ^{\prime}\right) ^{n} x\right) \\ & \leqq \left(\frac {| \lambda - t |}{t} C ^{-1} e\right) ^{n} \cdot q (x), \text{ whenever } 0 \leqq t / n \leqq 1. \end{aligned}
\]

Thus the right side of \eqref{eq:IX.10.2} is convergent for $|\arg \lambda| < \tan^{-1}(Ce^{-1})$, and so, by the sequential completeness of the space $X$, $T_{\lambda}x$ is well defined and weakly holomorphic in $\lambda$ for $|\arg \lambda| < \tan^{-1}(Ce^{-1})$. That is, for any $x \in X$ and $f \in X'$, the numerical function $f(T_t x)$ of $t$, $t > 0$, admits a holomorphic extension $f(T_{\lambda}x)$ for $|\arg \lambda| < \tan^{-1}(Ce^{-1})$; consequently, by the Hahn-Banach theorem, we see that $T_{\lambda}x$ is an extension of $T_t x$ for $|\arg \lambda| < \tan^{-1}(Ce^{-1})$. Next put $S_t = e^{-t} T_t$. Then $S_t' = e^{-t} T_t' - e^{-t} T_t$ and so, by $0 \leqq te^{-t} \leqq 1$ ( $0 \leqq t$) and (I), we easily see that the family of operators $\{(2^{-k} Ct S_t')^n\}$ is equi-continuous in $t > 0$ and $n \geqq 0$, in virtue of the equi-continuity of $\{T_t\}$ in $t > 0$. The equi-continuous semi-group $S_t$ of class $(C_0)$ satisfies the condition that $t > 0$ implies $S_t X \subseteq D(A - I) = D(A)$, where $(A - I)$ is the infinitesimal generator of $S_t$. Therefore, by the same reasoning applied above to $T_t$, we can prove that the weakly holomorphic extension $e^{-\lambda} T_{\lambda}$ of $S_t = e^{-t} T_t$ satisfies the estimate \eqref{eq:IX.10.3}. 
\end{proof}

By the way, we can prove the following 

\begin{corollaryn}[due to E. \textsc{Hille}] If, in particular, $X$ is a complex $B$-space and $\varlimsup\limits_{t\downarrow 0}\|tT_t'| < e^{-1}$, then $X = D(A)$.

\end{corollaryn}
\begin{proof}
 For a fixed $t > 0$, we have $\varlimsup\limits_{n \to \infty} \|(t/n) T_{t/n}'\| < e^{-1}$, and so the series 
\[
\sum\limits_{n=0} ^{\infty} (\lambda - t) ^{n} T _{t} ^{(n)} x / n! = \sum\limits_{n=0} ^{\infty} \frac {(\lambda - t) ^{n}}{t ^{n}} \frac {n ^{n}}{n !} \left(\frac {t}{n} T _{t/n} ^{\prime}\right) ^{n} x
\]
strongly converges in some sector 
\[
\{\lambda ; | \lambda - t | / t < 1 + \delta \quad \text{ with some } \quad \delta > 0 \}
\]
of the complex $\lambda$-plane. This sector surely contains $\lambda = 0$ in its interior. 


The implication (II) $\rightarrow$ (III). We have, by \eqref{eq:IX.4.10} in Chapter~\ref{chap:IX}, \ref{sec:IX.4}, 
\begin{equation}\label{eq:IX.10.4}
(\lambda R (\lambda ; A)) ^{n+1} x = \frac {\lambda^{n+1}}{n !} \int_{0} ^{\infty} e ^{- \lambda t} t ^{n} T _{t} x \dd t \quad \text{ for } \quad \operatorname{Re} (\lambda) > 0, x \in X.
\end{equation}

Hence, putting $S_{t} = e^{-t}T_{t}$, we obtain 
\[
\begin{aligned}
&((\sigma + 1 + i \tau) R (\sigma + 1 + i \tau ; A)) ^{n+1} x = \frac {(\sigma + 1 + i \tau) ^{n+1}}{n !} \int_{0} ^{\infty} e ^{-(\sigma + i \tau) t} t ^{n} S _{t} x \dd t, \sigma > 0.
\end{aligned}
\]

Let $\tau < 0$. Since the integrand is weakly holomorphic, we can deform, by the estimate \eqref{eq:IX.10.3} and Cauchy's integral theorem, the path of integration: $0 \leqq t < \infty$ to the ray: $re^{i\theta} (0 \leqq r < \infty)$ contained in the sector $0 < \arg \lambda < \tan^{-1}(2^{-k} Ce^{-1})$ of the complex $\lambda$-plane. We thus obtain 
\[
\begin{aligned}
&\quad ((\sigma + 1 + i \tau) R (\sigma + 1 + i \tau ; A)) ^{n+1} x = \frac {(\sigma + 1 + i \tau) ^{n+1}}{n !} \int_{0} ^{\infty} e ^ {- (\sigma + i \tau) r e ^{i \theta}} r ^{n} e ^{i n \theta} S _ {r e ^{i \theta}} x e ^{i \theta} \dd r,
\end{aligned}
\]
and so, by \eqref{eq:IX.10.3}, 
\[
\begin{aligned} & p \left(\left((\sigma + 1 + i \tau) R (\sigma + 1 + i \tau ; A)\right) ^{n+1} x\right) \\ & \leqq \sup\limits _{0 < r < \infty} p (S _ {r e ^{i \theta}} x) \frac {| (\sigma + 1 + i \tau) | ^{n+1}}{n !} \int_{0} ^{\infty} e ^{(- \sigma \cos \theta + \tau \sin \theta) r} r ^{n} \dd r \end{aligned}
\]
$\leqq q'(x) \frac{|\sigma + 1 + i\tau|^{n+1}}{|\tau \sin \theta - \sigma \cos \theta|^{n+1}}$, where $q'$ is a continuous semi-norm on $X$. 

A similar estimate is also obtained for the case $\tau > 0$. Hence, combined with \eqref{eq:IX.10.7} in Section~\hyperref[sec:IX.4]{$4$}, we have proved (III). 

The implication (III) $\rightarrow$ (I). For any continuous semi-norm $p$ on $X$, there exists a continuous semi-norm $q$ on $X$ such that $p((C_1\lambda R(\lambda;A))^n x)\leqq q(x)$ whenever $\operatorname{Re}(\lambda)\geqq 1 + \varepsilon, \varepsilon >0$, and $n\geqq 0$. Hence, if $\operatorname{Re}(\lambda_0)\geqq 1 + \varepsilon$, we have 
\[
\phi \left((\lambda - \lambda_{0}) ^{n} R (\lambda_{0}; A) ^{n} x\right) \leqq \frac {| \lambda - \lambda_{0} | ^{n}}{(C _{1} | \lambda_{0} |) ^{n}} q (x) \quad (n = 0, 1, 2, \dots).
\]

Thus, if $|\lambda - \lambda_0| / C_1 |\lambda_0| < 1$, the resolvent $R(\lambda, A)$ exists and is given by 
\begin{align}
R (\lambda ; A) x& = \sum\limits_{n=0} ^{\infty} (\lambda_{0} - \lambda) ^{n} R (\lambda_{0}; A) ^{n+1} x \quad \text{ such that } \notag \\
p (R (\lambda ; A) x)& \leqq (1 - C _{1} ^{-1} | \lambda_{0} | ^{-1} | \lambda - \lambda_{0} |) ^{-1} q (R (\lambda_{0}; A) x). \notag
\end{align}

Therefore, by (III), there exists an angle $\theta_0$ with $\pi / 2 < \theta_0 < \pi$ such that $R(\lambda; A)$ exists and satisfies the estimate 
\begin{equation}\label{eq:IX.10.5}
p(R(\lambda; A)x) \leqq \frac{1}{|\lambda|} q'(x) \text{ with a continuous semi-norm } q' \text{ on } X
\end{equation}
in the sectors $\pi / 2 \leqq \arg \lambda \leqq \theta_0$ and $-\theta_0 \leqq \arg \lambda \leqq -\pi / 2$ and also for $\operatorname{Re}(\lambda) \geqq 0$, when $|\lambda|$ is sufficiently large. 

Hence the integral 
\begin{equation}\label{eq:IX.10.6}
\hat {T} _{t} x = (2 \pi i) ^{-1} \int_ {C _{3}} e ^{\lambda t} R (\lambda ; A) x \dd \lambda \quad (t > 0, x \in X)
\end{equation}
converges if we take the path of integration $C_2 = \lambda(\sigma), -\infty < \sigma < \infty$, in such a way that $\lim\limits_{|\sigma| \uparrow \infty} |\lambda(\sigma)| = \infty$ and, for some $\varepsilon > 0$, 
\[
\pi / 2 + \varepsilon \leqq \arg \lambda (\sigma) \leqq \theta_{0} \quad \text{ and } \quad - \theta_{0} \leqq \arg \lambda (\sigma) \leqq - \pi / 2 - \varepsilon
\]
when $\sigma \uparrow +\infty$ and $\sigma \downarrow -\infty$, respectively; for $|\sigma|$ not large, $\lambda (\sigma)$ lies in the right half plane of the complex $\lambda$-plane. 

We shall show that $\hat{T}_{t}$ coincides with the semi-group $T_{t}$ itself. We first show that $\lim\limits_{t\downarrow0}\hat{T}_{t}x=x$ for all $x\in D(A)$. Let $x_{0}$ be any element $\in D(A)$, and choose any complex number $\lambda_{0}$ to the right of the contour $C_{2}$ of integration, and denote $(\lambda_{0}I-A)x_{0}=y_{0}$. Then, by the resolvent equation, 
\[
\begin{aligned}
&\hat {T} _{t} x _{0} = \hat {T} _{t} R (\lambda_{0}; A) y _{0}
\\
&= (2 \pi i) ^{-1} \int_ {C _{2}} e ^{\lambda t} R (\lambda ; A) R (\lambda_{0}; A) y _{0} \dd \lambda
\\
&= (2 \pi i) ^{-1} \int_ {C _{2}} e ^{\lambda t} (\lambda_{0} - \lambda) ^{-1} R (\lambda ; A) y _{0} \dd \lambda - (2 \pi i) ^{-1} \int_ {C _{2}} e ^{\lambda t} (\lambda_{0} - \lambda) ^{-1} R (\lambda_{0}; A) y _{0} \dd \lambda .
\end{aligned}
\]

The second integral on the right is equal to zero, as may be seen by shifting the path of integration to the left. Hence 
\[
\hat {T} _{t} x _{0} = (2 \pi i) ^{-1} \int_ {C _{3}} e ^{\lambda t} (\lambda_{0} - \lambda) ^{-1} R (\lambda ; A) y _{0} \dd \lambda , y _{0} = (\lambda_{0} I - A) x _{0}.
\]

Because of the estimate \eqref{eq:IX.10.5}, the passage to the limit $t \downarrow 0$ under the integral sign is justified, and so 
\[
\lim\limits _{t \downarrow 0} \hat {T} _{t} x _{0} = (2 \pi i) ^{-1} \int_ {C _{*}} (\lambda_{0} - \lambda) ^{-1} R (\lambda ; A) y _{0} \dd \lambda , y _{0} = (\lambda_{0} I - A) x _{0}.
\]

To evaluate the right hand integral, we make a closed contour out of the original path of integration $C_2$ by adjoining the arc of the circle $|\lambda| = r$ which is to the right of the path $C_2$, and throwing away that portion of the original path $C_2$ which lies outside the circle $|\lambda| = r$. The value of the integral along the new arc and the discarded arc tends to zero as $r \uparrow \infty$, in virtue of \eqref{eq:IX.10.5}. Hence the value of the integral is equal to the residue inside the new closed contour, that is, the value $R(\lambda_0; A) y_0 = x_0$. We have thus proved $\lim\limits_{t \downarrow 0} \hat{T}_t x_0 = x_0$ when $x_0 \in D(A)$. 

We next show that $\hat{T}'x = A\hat{T}_t x$ for $t > 0$ and $x \in X$. We have $R(\lambda; A)X = D(A)$ and $AR(\lambda; A) = \lambda R(\lambda; A) - I$, so that, by the
convergence factor $e^{\lambda t}$, the integral $(2\pi i)^{-1} \int_{C_2} e^{\lambda t} A R(\lambda; A)x \, \dd\lambda$ has a sense. This integral is equal to $A \hat{T}_t x$, as may be seen by approximating the integral \eqref{eq:IX.10.6} by Riemann sums and using the fact that $A$ is closed: $\lim\limits_{n \to \infty} x_n = x$ and $\lim\limits_{n \to \infty} A x_n = y$ imply $x \in D(A)$ and $Ax = y$. Therefore 
\[
A \hat {T} _{t} x = (2 \pi i) ^{-1} \int_ {C _{3}} e ^{\mu} A R (\lambda ; A) x \dd \lambda , t > 0.
\]

On the other hand, by differentiating \eqref{eq:IX.10.6} under the integral sign, we obtain 
\begin{equation}\label{eq:IX.10.7}
\hat {T} _{t} ^{\prime} x = (2 \pi i) ^{-1} \int_ {C _{2}} e ^{\lambda t} \lambda R (\lambda ; A) x \dd \lambda , t > 0.
\end{equation}

The difference of these two integrals is $(2\pi i)^{-1}\int_{C_{s}}e^{\lambda t}xd\lambda$, and the value of the last integral is zero, as may be seen by shifting the path of integration to the left. 

Thus we have proved that $\hat{x}(t) = \hat{T}_t x_0$, $x_0 \in D(A)$, satisfies \begin{enumerate}[i)]
\item $\lim\limits_{t \downarrow 0} \hat{x}(t) = x_0$,
\item $\dd\hat{x}(t)/\dd t = A\hat{x}(t)$ for $t > 0$
\item $\{\hat{x}(t)\}$ is by \eqref{eq:IX.10.6} of exponential growth when $t \uparrow \infty$.
\end{enumerate}
 On the other hand, since $x_0 \in D(A)$ and since $\{T_t\}$ is equi-continuous in $t \geqq 0$, we see that $x(t) = T_t x_0$ also satisfies $\lim\limits_{t \downarrow 0} x(t) = x_0$, $\dd x(t)/\dd t = Ax(t)$ for $t \geqq 0$, and $\{x(t)\}$ is bounded when $t \geqq 0$. Let us put $\hat{x}(t) - x(t) = y(t)$. Then $\lim\limits_{t \downarrow 0} y(t) = 0$, $\dd y(t)/\dd t = Ay(t)$ for $t > 0$ and $\{y(t)\}$ is of exponential growth when $t \uparrow \infty$. Hence we may consider the Laplace transform 
\[
L (\lambda ; y) = \int_{0} ^{\infty} e ^{- \lambda t} y (t) \dd t \text{ for large positive } \operatorname{Re} (\lambda).
\]

We have 
\[
\int_{\alpha} ^{\beta} e ^{- \lambda t} y ^{\prime} (t) \dd t = \int_{\alpha} ^{\beta} e ^{- \lambda t} A y (t) \dd t = A \int_{\alpha} ^{\beta} e ^{- \lambda t} y (t) \dd t, 0 \leqq \alpha < \beta < \infty ,
\]
by approximating the integral by Riemann sums and using the fact that $A$ is closed. By partial integration, we obtain 
\[
\int_{\alpha} ^{\beta} e ^{- \lambda t} y ^{\prime} (t) \dd t = e ^{- \lambda \beta} y (\beta) - e ^{- \lambda \alpha} y (\alpha) + \lambda \int_{\alpha} ^{\beta} e ^{- \lambda t} y (t) \dd t,
\]
which tends to $\lambda L(\lambda; y)$ as $\alpha \downarrow 0$, $\beta \uparrow \infty$. For, $y(0) = 0$ and $\{y(\beta)\}$ is of exponential growth as $\beta \uparrow \infty$. Thus again, by using the closure property of $A$, we obtain 
\[
A L (\lambda ; y) = \lambda L (\lambda ; y) \quad \text{ for large positive } \quad \operatorname{Re} (\lambda).
\]

Since the inverse $(\lambda I - A)^{-1}$ exists for $\operatorname{Re}(\lambda) > 0$, we must have $L(\lambda; y) = 0$ when $\operatorname{Re}(\lambda)$ is large positive. Thus, for any continuous linear functional 
$f\in X^{\prime}$, we have 
\[
\int_{0} ^{\infty} e ^{- \lambda t} f (y (t)) \dd t = 0 \quad \text{ when } \quad \operatorname{Re} (\lambda) \text{ is large positive. }
\]

We set $\lambda = \sigma + i\tau$ and put 
\[
g _{\sigma} (t) = e ^{- \sigma t} f (y (t)) \text{ or } = 0 \quad \text{ according as } \quad t \geqq 0 \quad \text{ or } \quad t < 0.
\]

Then, the above equality shows that the Fourier transform $(2\pi)^{-1}\int_{-\infty}^{\infty}e^{-i\tau t}g_{\sigma}(t)\dd t$ vanishes identically in $\tau$, $-\infty < \tau < \infty$, so that, by Fourier's integral theorem, $g_{\sigma}(t) = 0$ identically. Thus $f(y(t)) = 0$ and so we must have $y(t) = 0$ identically, in virtue of the Hahn-Banach theorem. 

Therefore $\hat{T}_t x = T_t x$ for all $t > 0$ and $x \in D(A)$. $D(A)$ being dense in $X$ and $\hat{T}_t, T_t$ both belonging to $L(X, X)$, we easily conclude that $\hat{T}_t x = T_t x$ for all $x \in X$ and $t > 0$. Hence, by defining $\hat{T}_0 = I$, we have $\hat{T}_t = T_t$ for all $t \geqq 0$. Hence, by \eqref{eq:IX.10.7}, $T_t' x = (2\pi i)^{-1} \int_{C_1} e^{\lambda t} \lambda R(\lambda; A) x \dd\lambda$, $t > 0$, and so, by \eqref{eq:IX.10.1} and \eqref{eq:IX.10.5}, we obtain 
\[
(T _{t/n} ^{\prime}) ^{n} x = T _{t} ^{(n)} x = (2 \pi i) ^{-1} \int_ {C _{3}} e ^{\lambda t} \lambda^{n} R (\lambda ; A) x \dd \lambda , t > 0.
\]

Hence 
\[
(t T _{t} ^{\prime}) ^{n}. x = (2 \pi i) ^{-1} \int_ {C _{s}} e ^{n \lambda t} (t \lambda) ^{n} R (\lambda ; A) x \dd \lambda , t > 0.
\]

Therefore, by (III), 
\[
\phi \left((t T _{t} ^{\prime}) ^{n} x\right) \leqq (2 \pi) ^{-1} q (x) \int_ {C _{s}} | e ^{n \lambda t} | t ^{n} | \lambda | ^{n-1} d | \lambda |.
\]

If $0 < t \leqq 1$, then the last integral is majorized by $C_{3}^{n}$ where $C_{3}$ is a certain positive constant. This we see, by splitting the integration path $C_{2}$ into the sum of that in the right half-plane $\operatorname{Re}(\lambda) \geqq 0$ and those in the left half-plane $\operatorname{Re}(\lambda) < 0$, and remembering the integral representation of the $\Gamma$-function.
\end{proof}

References. The result of the present section is due to K. \textsc{Yosida} \cite{yosida349}. See also E. \textsc{Hille} \cite{hille99} and E. \textsc{Hille}-R. S. \textsc{Phillips} \cite{hille97}. 


\yossection{11}{Fractional Powers of Closed Operators}
\label{sec:IX.11}

Let $X$ be a $B$-space, and $\{T_t; t \geqq 0\} \subseteq L(X, X)$ an equi-continuous semi-group of class $(C_0)$. We introduce 
\begin{equation}\label{eq:IX.11.1}
\begin{aligned}
&f _{t, \alpha} (\lambda)  = \frac {1}{2 \pi i} \int_{\sigma - i \infty} ^{\sigma + i \infty} e ^ {i \lambda - t s ^{\alpha}} \dd z (\sigma > 0, t > 0, \lambda \geqq 0, 0 < \alpha < 1), = 0 \quad (\text{ when } \lambda < 0),
\end{aligned}
\end{equation}
where the branch of $z^{\alpha}$ is so taken that $\operatorname{Re}(z^{\alpha}) > 0$ for $\operatorname{Re}(z) > 0$. This branch is a one-valued function in the $z$-plane cut along the negative real axis. The convergence of the integral \eqref{eq:IX.11.1} is apparent in virtue of the convergence factor $e^{-tz^{\alpha}}$. Following S. \textsc{Bochner} \cite{bochner21} and R. S. \textsc{Phillips} \cite{phillips255}, we can show that the operators 
\begin{equation}\label{eq:IX.11.2}
\begin{aligned} \hat {T} _{t, \alpha} x & = \hat {T} _{t} x = \int_{0} ^{\infty} f _{t, \alpha} (s) T _{s} x \dd s (t > 0), \\ & = x (t = 0), \end{aligned}
\end{equation}
constitute an equi-continuous semi-group of class $(C_{0})$. Moreover, we can show that $\{\hat{T}_{t}\}$ is a holomorphic semi-group (K. \textsc{Yosida} \cite{yosida351} and V. \textsc{Balakrishnan} \cite{balakrishnan8}). The infinitesimal generator $\hat{A} = \hat{A}_{\alpha}$ of $\hat{T}_{t}$ is connected with the infinitesimal generator $A$ of $T_{t}$ by 
\begin{equation}\label{eq:IX.11.3}
\hat {A} _{\alpha} x = - (- A) ^{\alpha} x \quad \text{ for } \quad x \in D (A),
\end{equation}
where the fractional power $(-A)^{\alpha}$ of $(-A)$ is given by 
\begin{equation}\label{eq:IX.11.4}
(- A) ^{\alpha} x = \frac {\sin \alpha \pi}{\pi} \int_{0} ^{\infty} \lambda^{\alpha -1} (\lambda I - A) ^{-1} (- A x) \dd \lambda \quad \text{ for } \quad x \in D (A)
\end{equation}
and also by 
\begin{equation}\label{eq:IX.11.5}
(- A) ^{\alpha} x = \Gamma (- \alpha) ^{-1} \int_{0} ^{\infty} \lambda^{- \alpha -1} (T _{\lambda} - I) x \dd \lambda \quad \text{ for } \quad x \in D (A).
\end{equation}

Formula \eqref{eq:IX.11.4} and \eqref{eq:IX.11.5} were obtained by V. \textsc{Balakrishnan}. For the resolvent of $\hat{A}_{\alpha}$, we have the following formula due to T. \textsc{Kato}: 
\begin{equation}\label{eq:IX.11.6}
(\mu I - \hat {A} _{\alpha}) ^{-1} = \frac {\sin \alpha \pi}{\pi} \int_{0} ^{\infty} (r I - A) ^{-1} \frac {r ^{\alpha}}{\mu^{2} - 2 \mu r ^{\alpha} \cos \alpha \pi + r ^{2 \alpha}} \dd r.
\end{equation}

In this way, we see the abundance of holomorphic semi-groups among the class of equi-continuous semi-groups of class $(C_{0})$. 

To prove the above result, we have to investigate the properties of the function $f_{t,\alpha}(\lambda)$ in a series of propositions. 

\begin{proposition}\label{prop:IX.11.Proposition.1} We have 
\begin{equation}\label{eq:IX.11.7}
e ^ {- t a ^{\alpha}} = \int_{0} ^{\infty} e ^{- \lambda a} f _{t, \alpha} (\lambda) \dd \lambda (t > 0, a > 0).
\end{equation}
\end{proposition}

\begin{proof}
 It is easy to see, by the convergence factor $e^{-z^{\alpha}t}$, that the function $f_{t,\alpha}(\lambda)$ is of exponential growth in $\lambda$. By Cauchy's integral theorem, integral \eqref{eq:IX.11.1} is independent of $\sigma > 0$. Let $a > \sigma = \operatorname{Re}(z) > 0$. Then, by Cauchy's theorem of residue, 
\[ \begin{aligned} \int_{0} ^{\infty} e ^{- \lambda a} f _{t, \alpha} (\lambda) \dd \lambda & = \frac {1}{2 \pi i} \int_{\sigma - i \infty} ^{\sigma + i \infty} \left[ \frac {e ^{\lambda (z - a)}}{z - a} \right] _{\lambda = 0} ^{\lambda = \infty} e ^ {- z ^{\alpha} t} \dd z \\ & = \frac {- 1}{2 \pi i} \int_{\sigma - i \infty} ^{\sigma + i \infty} \frac {1}{z - a} e ^ {- z ^{\alpha} t} \dd z = e ^ {- t a ^{\alpha}}. \end{aligned}
\]
\end{proof}

\begin{proposition}\label{prop:IX.11.Proposition.2} We have 
\begin{equation}\label{eq:IX.11.8}
f _{t, \alpha} (\lambda) \geqq 0 \quad \text{ for all } \quad \lambda > 0.
\end{equation}
\end{proposition}

\begin{proof}
 If we set $a^{\alpha} = g(a)$, $e^{-tx} = h(x)$, then 
\[
\begin{aligned}
(-1)^{n-1} g^{(n)}(a) &\geqq 0 \quad (n = 1, 2, \ldots), \quad g(a) \geqq 0, \\
\text{and } (-1)^{n} h^{(n)}(x) &\geqq 0 \quad (n = 0, 1, \ldots), \quad \text{when } a \geqq 0 \text{ and } x \geqq 0.
\end{aligned}
\]

Hence $k(a) = h(g(a)) = e^{-ta^{\alpha}}$ satisfies 
\begin{equation}\label{eq:IX.11.9}
\begin{yoslines}& (- 1) ^{n} k ^{(n)} (a) = (- 1) h ^{\prime} (x) (- 1) ^{n-1} g ^{(n)} (a)\\& + \sum\limits_{(p)} C _ {(p _{0} p _{1} \dots p _{v})} ^{(n)} (- 1) ^ {p _{0}} h ^ {(p _{0})} (x) (- 1) ^ {p _{1}} g ^ {(p _{1} + 1)} (a) \dots (- 1) ^ {p _{v}} g ^ {(p _{v} + 1)} (a)\\& \left(C _{(p)} ^{(n)} \geqq 0, p _{0} \geqq 2, p _{1} \geqq 0, \dots , p _{v} \geqq 0 \text{ with } p _{0} \leqq \sum\limits_{i=1} ^{v} p _{i} = n, \text{ and } v \text{ arbitrary }\right)\\& \geqq 0, (n = 0, 1, \dots).\end{yoslines}
\end{equation}

That is, the function $k(a) = e^{-ta^{\alpha}}$ is completely monotone in $a \geqq 0$. 

We next prove the Post-Widder inversion formula 
\begin{equation}\label{eq:IX.11.10}
f _{t, \alpha} (\lambda) = \lim\limits _{n \rightarrow \infty} \frac {(- 1) ^{n}}{n !} \left(\frac {n}{\lambda}\right) ^{n+1} k ^{(n)} \left(\frac {n}{\lambda}\right), \quad \lambda > 0
\end{equation}
so that, by \eqref{eq:IX.11.9}, $f_{t,\alpha}(\lambda) \geqq 0$. The proof of \eqref{eq:IX.11.10} is obtained as follows. We find, differentiating \eqref{eq:IX.11.7} $n$-times, 
\[
k ^{(n)} \left(\frac {n}{\lambda}\right) = (- 1) ^{n} \int_{0} ^{\infty} s ^{n} e ^{- s n / \lambda} f _{t, \alpha} (s) \dd s.
\]

Substituting this into the right side of \eqref{eq:IX.11.10}, we get 
\[
\lim\limits _{n \rightarrow \infty} \frac {n ^{n+1}}{e ^{n} n !} \frac {1}{\lambda} \int_{0} ^{\infty} \left[ \frac {s}{\lambda} \cdot \exp \left(1 - \frac {s}{\lambda}\right)\right] ^{n} f _{t, \alpha} (s) \dd s.
\]

Since 
\[
\lim\limits _{n \rightarrow \infty} n ^{n} / \sqrt {2 \pi n} e ^{n} n! = 1 \quad (\text{ Stirling's formula }),
\]
we have to prove that 
\begin{equation}\label{eq:IX.11.11}
f _{t, \alpha} (\lambda_{0}) = \lim\limits _{n \rightarrow \infty} \frac {\sqrt {2 \pi}}{\lambda_{0}} \int_{0} ^{\infty} n ^{3/2} \left[ \frac {s}{\lambda_{0}} \cdot \exp \left(1 - \frac {s}{\lambda_{0}}\right)\right] ^{n} f _{t, \alpha} (s) \dd s, \lambda_{0} > 0.
\end{equation}

Let $\eta$ be a fixed positive number such that $\eta < \lambda_0$. We decompose the last integral into three parts, 
\[
\int_{0} ^{\infty} = \int_{0} ^ {\lambda_{0} - \eta} + \int_ {\lambda_{0} - \eta} ^ {\lambda_{0} + \eta} + \int_ {\lambda_{0} + \eta} ^{\infty} = J _{1} + J _{2} + J _{3}.
\]

Since $x \cdot \exp(1 - x)$ increases monotonically in [$0$, $1$] from $0$ to $1$, we see, by the boundedness of $f_{t,\alpha}(s)$ in $s$, that $\lim\limits_{n\to \infty}J_1 = 0$. Next, since $x\cdot \exp (1 - x)$ decreases monotonically in $[1,\infty]$ from $1$ to $0$, we have 
\[
\frac {\lambda_{0} + \eta}{\lambda_{0}} \cdot \exp \left(1 - \frac {\lambda_{0} + \eta}{\lambda_{0}}\right) < \beta < 1,
\]
and so, since $f_{t,\alpha}(s)$ is of exponential growth as $s \uparrow \infty$, 
\[
\left| J _{3} \right| \leqq n ^{3/2} e ^ {n _{0}} \beta^ {n - n _{0}} \int_{0} ^{\infty} \left(\frac {s}{\lambda_{0}}\right) ^ {n _{0}} \exp \left(- \frac {n _{0} s}{\lambda_{0}}\right)\left| f _{t, \alpha} (s) \right| \dd s \rightarrow 0 \text{ as } n \rightarrow \infty .
\]

By the continuity of $f_{t,\alpha}(s)$ in $s$, we have, for any positive number $\varepsilon$, $f_{t,\alpha}(\lambda_0) - \varepsilon \leqq f_{t,\alpha}(s) \leqq f_{t,\alpha}(\lambda_0) + \varepsilon$ whenever $\lambda_0 - \eta \leqq s \leqq \lambda_0 + \eta$, if we take $\eta > 0$ sufficiently small. Thus 
\begin{equation}\label{eq:IX.11.12}
\left(f _{t, \alpha} (\lambda_{0}) - \varepsilon\right) J _{0} \leqq J _{2} \leqq \left(f _{t, \alpha} (\lambda_{0}) + \varepsilon\right) J _{0},
\end{equation}
where 
\begin{equation}\label{eq:IX.11.13}
J _{0} = \int_ {\lambda_{0} - \eta} ^ {\lambda_{0} + \eta} n ^{3/2} \left[ \frac {s}{\lambda_{0}} \cdot \exp \left(1 - \frac {s}{\lambda_{0}}\right) \right] ^{n} \dd s.
\end{equation}

The whole preceding argument is true for the particular case of the completely monotone function 
\[
k (a) = a ^{-1} = \int_{0} ^{\infty} e ^{- \lambda a} \dd \lambda .
\]

In this case, $k^{(n)}(n / \lambda_0) = (-1)^n n!(\lambda_0 / n)^{n+1}$. Substituting this in \eqref{eq:IX.11.10}, we find that \eqref{eq:IX.11.10} holds for $f_{t,\alpha}(\lambda) \equiv 1$. Since \eqref{eq:IX.11.10} and \eqref{eq:IX.11.11} are equivalent, \eqref{eq:IX.11.11} must also hold for $f_{t,\alpha}(\lambda) \equiv 1$. Thus, since $\lim\limits_{n\to \infty}J_1 = 0$ and $\lim\limits_{n\to \infty}J_3 = 0$ for a general $f_{t,\alpha}$, we obtain 
\[
1 = \lim\limits _{n \rightarrow \infty} \sqrt {2 \pi} \lambda_{0} ^{-1} J _{0}.
\]

Therefore, by \eqref{eq:IX.11.12}, we get \eqref{eq:IX.11.11} and the equivalent formula \eqref{eq:IX.11.10} is proved. 
\end{proof}

\begin{proposition}\label{prop:IX.11.Proposition.3} 
\begin{align}
&\int_{0} ^{\infty} f _{t, \alpha} (\lambda) \dd \lambda = 1, \label{eq:IX.11.14} \\
f _{t + s, \alpha} (\lambda)& = \int_{0} ^{\infty} f _{t, \alpha} (\lambda - \mu) f _{s, \alpha} (\mu) \dd \mu . \label{eq:IX.11.15}
\end{align}
\end{proposition}

\begin{proof}
 Since the function $f_{t,\alpha}(\lambda)$ is non-negative, we have, by the Lebesgue-Fatou lemma and \eqref{eq:IX.11.7}, 
\[
\int_{0} ^{\infty} \varliminf\limits_{a \downarrow 0} \left(e ^{- \lambda a} f _{t, \alpha} (\lambda)\right) \dd \lambda \leqq \varliminf\limits_{a \downarrow 0} e ^ {- t a ^{\alpha}} = 1. \]

Thus $f_{t,\alpha}(\lambda)$ is integrable with respect to $\lambda$ over $(0,\infty)$ and so, again by 
the Lebesgue-Fatou lemma and \eqref{eq:IX.11.7}, we obtain \eqref{eq:IX.11.14}. We next have, by \eqref{eq:IX.11.7}, 
\[
\begin{aligned}
&\int e ^{- \lambda a} \left\{\int f _{t, \alpha} (\lambda - \mu) f _{s, \alpha} (\mu) \dd \mu \right\} \dd \lambda = \int e ^{-(\lambda - \mu) a} f _{t, \alpha} (\lambda - \mu) d (\lambda - \mu) \cdot \int e ^{- \mu a} f _{s, \alpha} (\mu) \dd \mu
\\
&= e ^ {- t a ^{\alpha}} e ^ {- s a ^{\alpha}} = e ^ {- (t + s) a ^{\alpha}} = \int e ^{- \lambda a} f _{t + s, \alpha} (\lambda) \dd \lambda , a > 0.
\end{aligned}
\]

Hence, by inverting the Laplace transform as in the preceding section, we get \eqref{eq:IX.11.15}. 
\end{proof}

\begin{proposition}\label{prop:IX.11.Proposition.4} We have 
\begin{equation}\label{eq:IX.11.16}
\int_{0} ^{\infty} \partial f _{t, \alpha} (\lambda) / \partial t \cdot \dd \lambda = 0, \quad t > 0.
\end{equation}
\end{proposition}

\begin{proof}
 By deforming the path of integration in \eqref{eq:IX.11.1} to the union of two paths $re^{-i\theta}(-\infty<-r<0)$ and $re^{i\theta}(0<r<\infty)$, where $\pi/2\leqq\theta\leqq\pi$, we obtain 
\begin{equation}\label{eq:IX.11.17}
\begin{aligned}
&f _{t, \alpha} (s) = \frac {1}{\pi} \int_{0} ^{\infty} \exp (s r \cdot \cos \theta - t r ^{\alpha} \cos \alpha \theta) \times \sin (s r \cdot \sin \theta - t r ^{\alpha} \sin \alpha \theta + \theta) \dd r.
\end{aligned}
\end{equation}

Similarly, by deforming the path of integration in 
\[
\partial f _{t, \alpha} (\lambda) / \partial t = \frac {1}{2 \pi i} \int_{\sigma - i \infty} ^{\sigma + i \infty} e ^ {z \lambda - z ^{\alpha} t} (- z ^{\alpha}) \dd z
\]
to the union of two paths, $re^{-i\theta}(-\infty < -r < 0)$ and $re^{i\theta}(0 < r < \infty)$, we obtain 
\begin{equation}\label{eq:IX.11.18}
\begin{aligned} f _{t, \alpha} ^{\prime} (s) & = \partial f _{t, \alpha} (s) / \partial t = \frac {(- 1)}{\pi} \int_{0} ^{\infty} \exp (s r \cdot \cos \theta - t r ^{\alpha} \cos \alpha \theta) \\ & \times \sin (s r \cdot \sin \theta - t r ^{\alpha} \sin \alpha \theta + \alpha \theta + \theta) r ^{\alpha} \dd r. \end{aligned}
\end{equation}

If we take 
\[
\theta = \theta_{\alpha} = \pi / (1 + \alpha),
\]
then 
\begin{equation}\label{eq:IX.11.19}
f _{t, \alpha} ^{\prime} (s) = \frac {1}{\pi} \int_{0} ^{\infty} \exp \left((s r + t r ^{\alpha}) \cos \theta_{\alpha}\right) \cdot \sin \left((s r - t r ^{\alpha}) \sin \theta_{\alpha}\right) r ^{\alpha} \dd r.
\end{equation}

Thus we see, by the factor $r^{\alpha}$ ( $0 < \alpha < 1$), that $f_{t,\alpha}'(s)$ is integrable with respect to $s$ over $(0, \infty)$. Hence, by differentiating \eqref{eq:IX.11.14} with respect to $t$, we get \eqref{eq:IX.11.16}. 
\end{proof}

We are now able to prove 

\begin{theorem}\label{thm:IX.11.Theorem.1} $\{\hat{T}_t\}$ is a holomorphic semi-group. 
\end{theorem}

\begin{proof}
 That $\{\hat{T}_t\}$ enjoys the semi-group property $\hat{T}_t\hat{T}_s = \hat{T}_{t+s}$ ( $t, s > 0$) is clear from \eqref{eq:IX.11.2} and \eqref{eq:IX.11.15}. We have, by \eqref{eq:IX.11.2} and \eqref{eq:IX.11.17} with $\theta = \theta_{\alpha}$, 
\begin{equation}\label{eq:IX.11.20}
\begin{aligned}
&\hat {T} _{t} x = \frac {1}{\pi} \int_{0} ^{\infty} T _{s} x \dd s \int_{0} ^{\infty} \exp ((s r + t r ^{\alpha}) \cos \theta_{\alpha}) \times \sin ((s r - t r ^{\alpha}) \sin \theta_{\alpha} + \theta_{\alpha}) \dd r,
\end{aligned}
\end{equation}
which gives by the change of variables 
\begin{gather}s = v t ^{1 / \alpha}, r = u t ^{-1 / \alpha}, \hat {T} _{t} x= \frac {1}{\pi} \int_{0} ^{\infty} T _ {v t ^{1 / \alpha}} \cdot x \dd v \int_{0} ^{\infty} \exp \left((u v + u ^{\alpha}) \cos \theta_{\alpha}\right) \label{eq:IX.11.21}
\\\quad \times \sin \left((u v - u ^{\alpha}) \sin \theta_{\alpha} + \theta_{\alpha}\right) \dd u.  \tag{20'} \label{eq:IX.11.20'}
\end{gather}

The second integral on the right is exactly $\pi \cdot f_{1,\alpha}(v)$, and so, by the equi-boundedness of $\{\|T_t x\|\}$ in $t \geqq 0$, we see, by \eqref{eq:IX.11.14}, that 
\begin{equation}\label{eq:IX.11.22}
\left\| \hat {T} _{t} x \right\| \leqq \sup\limits _{t \geqq 0} \left\| T _{t} x \right\| \int_{0} ^{\infty} f _{1, \alpha} (v) \dd v = \sup\limits _{t \geqq 0} \left\| T _{t} x \right\|.
\end{equation}

By passing to the limit $t \downarrow 0$ in $\eqref{eq:IX.11.20}$, which is justified by the integrability of $f_{1,\alpha}(v)$ over $[0, \infty)$, we obtain, by \eqref{eq:IX.11.14}, 
\[
\slim _{t \downarrow 0} \hat {T} _{t} x = \int_{0} ^{\infty} f _{1, \alpha} (v) \dd v \cdot x = x.
\]

Hence $\{\hat{T}_t\}$ is an equi-continuous semi-group of class $(C_0)$ such that \eqref{eq:IX.11.22} holds. 


By the integrability of $f_{t,\alpha}'(s) = \partial f_{t,\alpha}(s) / \partial t$ over $[0, \infty)$ and the equicontinuity of $\{T_t\}$, we obtain, by differentiating \eqref{eq:IX.11.2} with respect to $t$ under the integral sign, 
\begin{equation}\label{eq:IX.11.23}
\begin{aligned}
\hat {T} _{t} ^{\prime} x &= \int_{0} ^{\infty} f _{t, \alpha} ^{\prime} (s) T _{s} x \dd s
\\
&= \frac {1}{\pi} \int_{0} ^{\infty} T _{s} x \dd s \int_{0} ^{\infty} \exp ((s r + t r ^{\alpha}) \cos \theta_{\alpha}) \cdot \sin ((s r - t r ^{\alpha}) \sin \theta_{\alpha}) r ^{\alpha} \dd r,
\end{aligned}
\end{equation}
which is, by the change of variables \eqref{eq:IX.11.21}, 
\[
= \int_{0} ^{\infty} \left(T _ {v t ^{1 / \alpha}} \cdot x\right) \cdot f _{1, \alpha} ^{\prime} (v) \dd v \cdot t ^{-1}.
\]

Thus, by the integrability of $f_{1,\alpha}'(v)$ over $[0, \infty)$ and the equi-continuity of $\{T_t\}$ in $t \geqq 0$, we see that 
\[
\varlimsup\limits_{t \downarrow 0} \| t \hat {T} _{t} ^{\prime} \| < \infty ,
\]
that is, $\{\hat{T}_t\}$ is a holomorphic semi-group.
\end{proof}

\begin{theorem}\label{thm:IX.11.Theorem.2} The infinitesimal generator $\hat{A}_{\alpha}$ of $\hat{T}_t$ is connected with the infinitesimal generator $A$ of $T$ by \eqref{eq:IX.11.3}, where $(-A)^{\alpha}$ is defined by \eqref{eq:IX.11.4} and also by \eqref{eq:IX.11.5}. We have, moreover, \eqref{eq:IX.11.6}. 
\end{theorem}

\begin{proof}
 By \eqref{eq:IX.11.16} and \eqref{eq:IX.11.23}, we obtain 
\begin{equation}\label{eq:IX.11.24}
\begin{aligned}
&\hat {T} _{t} ^{\prime} x = \frac {1}{\pi} \int_{0} ^{\infty} (T _{s} - I) x \dd s \int_{0} ^{\infty} \exp ((s r + t r ^{\alpha}) \cos \theta_{\alpha}) \times \sin ((s r - t r ^{\alpha}) \sin \theta_{\alpha}) r ^{\alpha} \dd r.
\end{aligned}
\end{equation}

If $x \in D(A)$, then $\slim\limits_{s \downarrow 0} s^{-1}(T_s - I)x = Ax$ and $\|(T_s - I) \cdot x\|$ is bounded in $s \geqq 0$. Thus, we obtain, letting $t \downarrow 0$ in \eqref{eq:IX.11.24}, 
\[
\begin{aligned} \slim _{t \downarrow 0} \hat {T} _{t} ^{\prime} x & = \frac {1}{\pi} \int_{0} ^{\infty} (T _{s} - I) x \dd s \int_{0} ^{\infty} \exp (s r \cdot \cos \theta_{\alpha}) \cdot \sin (s r \cdot \sin \theta_{\alpha}) r ^{\alpha} \dd r \\ & = (- \Gamma (- \alpha)) ^{-1} \int_{0} ^{\infty} s ^{- \alpha -1} (T _{s} - I) x \dd s, \end{aligned}
\]
because, by the $\Gamma$-function formulae 
\begin{equation}\label{eq:IX.11.25}
\Gamma (z) = c ^{z} \int_{0} ^{\infty} e ^{-cr} r ^{z-1} \dd r \left(\operatorname{Re} (z) > 0, \operatorname{Re} (c) > 0\right)
\end{equation}
and 
\begin{equation}\label{eq:IX.11.26}
\Gamma (z) \Gamma (1 - z) = \pi / \sin \pi z,
\end{equation}
we obtain, by $(\alpha + 1)\theta_{\alpha} = \pi$, 
\[
\begin{aligned}
&\frac {1}{\pi} \int_{0} ^{\infty} \exp (s r \cdot \cos \theta_{\alpha}) \cdot \sin (s r \cdot \sin \theta_{\alpha}) r ^{\alpha} \dd r = (\pi i) ^{-1} \operatorname{Im} \left\{\int_{0} ^{\infty} e ^ {- r (- s e ^ {i \theta_{\alpha}})} r ^{\alpha} \dd r \right\}
\\
&= (\pi i) ^{-1} \operatorname{Im} ((- s e ^ {i \theta_{\alpha}}) ^{- \alpha -1}) \Gamma (1 + \alpha) = s ^{- \alpha -1} \pi^{-1} \sin (\alpha \pi) \Gamma (1 + \alpha) = s ^{- \alpha -1} \frac {- \Gamma (1 + \alpha)}{\Gamma (- \alpha) \Gamma (1 + \alpha)}
\\
&= (- \Gamma (- \alpha)) ^{-1} s ^{- \alpha -1} .
\end{aligned}
\]

Thus, by $\hat{T}_t' x = \hat{A}_{\alpha} \hat{T}_t x$ (when $t > 0$), the continuity at $t = 0$ of $\hat{T}_t x$ and the closure property of the infinitesimal generator $\hat{A}_{\alpha}$, we obtain 
\[
\hat {A} _{\alpha} x = (- \Gamma (- \alpha)) ^{-1} \int_{0} ^{\infty} s ^{- \alpha -1} (T _{s} - I) x \dd s \quad \text{ when } \quad x \in D (A).
\]

Hence, by \eqref{eq:IX.11.25}, \eqref{eq:IX.11.26} and $(tI - A)^{-1} = \int_{0}^{\infty} e^{-ts} T_s \dd s$, we obtain 
\[
\begin{aligned} \hat {A} _{\alpha} x & = \Gamma (- \alpha) ^{-1} \Gamma (1 + \alpha) ^{-1} \int_{0} ^{\infty} \left\{\int_{0} ^{\infty} e ^{-st} t ^{\alpha} \dd t \right\} (I - T _{s}) x \dd s \\ & = \frac {\sin \alpha \pi}{\pi} \int_{0} ^{\infty} t ^{\alpha} ((t I - A) ^{-1} - t ^{-1} I) x \dd t \\ & = \frac {\sin \alpha \pi}{\pi} \int_{0} ^{\infty} t ^{\alpha -1} (t I - A) ^{-1} A x \dd t \quad \text{ for } \quad x \in D (A). \end{aligned}
\]

Finally we have, by taking $\theta = \pi$ in \eqref{eq:IX.11.17} and \eqref{eq:IX.11.2}, 
\[
\begin{gathered}
(\mu I - \hat {A} _{\alpha}) ^{-1} = \int_{0} ^{\infty} e ^{- \mu t} \hat {T} _{t} \dd t
\\
= \pi^{-1} \int_{0} ^{\infty} \dd r \int_{0} ^{\infty} e ^{-sr} T _{s} \dd s \int_{0} ^{\infty} \exp (- \mu t - t r ^{\alpha} \cos \alpha \pi) \cdot \sin (t r ^{\alpha} \sin \alpha \pi) \dd t
\\
= \pi^{-1} \int_{0} ^{\infty} (r I - A) ^{-1} \left\{\int_{0} ^{\infty} \exp (- \mu t - t r ^{\alpha} \cos \alpha \pi) \cdot \sin (t r ^{\alpha} \sin \alpha \pi) \dd t \right\} \dd r
\\
= \frac {\sin \alpha \pi}{\pi} \int_{0} ^{\infty} (r I - A) ^{-1} \frac {r ^{\alpha}}{\mu^{2} - 2 r ^{\alpha} \mu \cos \alpha \pi + r ^{2 \alpha}} \dd r.
\end{gathered}
\]
\end{proof}

\begin{remarku} Formula \eqref{eq:IX.11.2} was devised by S. \textsc{Bochner} \cite{bochner21} without detailed proof. Cf. R. S. \textsc{Phillips} \cite{phillips255}. That $\hat{T}_t$ is a holomorphic semi-group was proved by K. \textsc{Yosida} \cite{yosida351}, V. \textsc{Balakrishnan} \cite{balakrishnan8} and T. \textsc{Kato} \cite{kato135}. Formulae \eqref{eq:IX.11.4} and \eqref{eq:IX.11.5} are due to V. \textsc{Balakrishnan} \cite{balakrishnan8}, who, by virtue of \eqref{eq:IX.11.4}, defined the fractional power $(-A)^{\alpha}$ of a closed linear operator $A$ satisfying the condition: 
\begin{equation}
\label{eq:IX.11.27}
\begin{minipage}{0.92\linewidth}
\everymath{}\let\limits\nolimits
the resolvent $R (\lambda; A) = (\lambda I - A) ^{-1}$ exists for $\operatorname{Re} (\lambda) > 0$ and $\sup\limits _{\operatorname{Re} (\lambda) > 0} | \operatorname{Re} (\lambda) | \cdot \| R (\lambda; A) \| < \infty.$
\end{minipage}
\end{equation}
\end{remarku}

He also proved that $(-A)^{a}$ enjoys properties to be demanded for the fractional power. In fact we have 

\begin{theorem}\label{thm:IX.11.Theorem.3} Let a closed linear operator $A$ satisfy condition \eqref{eq:IX.11.27}. Then, by \eqref{eq:IX.11.4}, a linear operator $(-A)^{\alpha}$ is defined, and we have 
\begin{align}
(- A) ^{\alpha} (- A) ^{\beta} x& = (- A) ^{\alpha + \beta} x \quad \text{ if } \quad x \in D (A ^{2}) \quad \text{ and } \quad 0 < \alpha , \beta \notag \\
\text{ with } \quad \alpha + \beta& < 1, \label{eq:IX.11.28} \\
\slim _{\alpha \uparrow 1} (- A) ^{\alpha} x& = - A x \quad \text{ if } \quad x \in D (A), \label{eq:IX.11.29} \\
\slim _{\alpha \downarrow 0} (- A) ^{\alpha} x& = x \quad \text{ if } \quad \slim _{\lambda \downarrow 0} \lambda R (\lambda ; A) x = 0, \label{eq:IX.11.30}
\end{align}
and, if $A$ is the infinitesimal generator of an equi-continuous semi-group $T_{t}$ of class $(C_0)$, 
\begin{equation}(A _{\alpha}) _{\beta} = A _{\alpha \beta}, \text{ where } A _{\alpha} \text{ is the operator } \hat {A} _{\alpha} \text{ defined through } \text{ Kato's formula } \eqref{eq:IX.11.6}. \label{eq:IX.11.31}\end{equation}
\end{theorem}

\begin{remarku} The last formula \eqref{eq:IX.11.31} is due to J. \textsc{Watanabe} \cite{watanabe331}. 
\end{remarku}

\begin{proof}
 $\left\|r^{\alpha-1}(rI-A)^{-1}(-Ax)\right\|$ is, by \eqref{eq:IX.11.27}, of order $O(r^{\alpha-2})$ when $r\uparrow\infty$, and, by $(rI-A)^{-1}(-Ax)=x-r(rI-A)^{-1}x$ and \eqref{eq:IX.11.27}, it is of order $O(r^{\alpha-1})$ when $r\downarrow0$. Thus the right side of \eqref{eq:IX.11.4} is convergent. 


It is clear that $x \in D(A^{2})$ implies $(-A)^{\beta}x \in D(A)$. For, by approximating the integral by Riemann sums and making use of the closure property of $A$, we have $(-A)^{\beta}x \subseteq D(A)$. We can thus define $(-A)^{\alpha}(-A)^{\beta}x$. 
\[
(- A) ^{\alpha} (- A) ^{\beta} x = \frac {\sin \alpha \pi}{\pi} \frac {\sin \beta \pi}{\pi} \int_{0} ^{\infty} \int_{0} ^{\infty} \lambda^{\beta -1} \mu^{\alpha -1} R (\lambda ; A) R (\mu ; A) A ^{2} x \dd \lambda \dd \mu
\]
can be rewritten, decomposing the domain of integration into those for which $\lambda \geqq \mu$ and $\lambda < \mu$, as follows. 
\[
\frac {\sin \alpha \pi}{\pi} \frac {\sin \beta \pi}{\pi} \int_{0} ^{1} \left(\sigma^{\beta -1} + \sigma^{\alpha -1}\right) \dd \sigma \int_{0} ^{\infty} \lambda^{\alpha + \beta -1} R (\lambda \sigma ; A) R (\lambda ; A) A ^{2} x \dd \lambda .
\]

By the resolvent equation $R(\lambda; A) - R(\mu; A) = (\mu - \lambda)R(\lambda; A)R(\mu; A)$ and $R(\lambda; A)(-A) = I - \lambda R(\lambda; A)$ valid on $D(A)$, we obtain 
$R(\lambda \sigma ;A)R(\lambda ;A)A^{2}x = (1 - \sigma)^{-1}\{-\sigma R(\lambda \sigma ;A) + R(\lambda ;A)\} (-Ax),$ and hence 
\[
\begin{aligned}
&(- A) ^{\alpha} (- A) ^{\beta} x = \frac {\sin \alpha \pi}{\pi} \frac {\sin \beta \pi}{\pi} \slim _{t \uparrow 1} \int_{0} ^{t} (\sigma^{\beta -1} + \sigma^{\alpha -1}) (1 - \sigma) ^{-1} \dd \sigma \\
&\times \int_{0} ^{\infty} \lambda^{\alpha + \beta -1} (- \sigma R (\lambda \sigma ; A) + R (\lambda ; A)) (- A x) \dd \lambda \\
&= \int_{0} ^{\infty} \left(\frac {\sin \alpha \pi}{\pi} \frac {\sin \beta \pi}{\pi} \int_{0} ^{1} \frac {\sigma^{\beta -1} + \sigma^{\alpha -1} - \sigma^{- \alpha} - \sigma^{- \beta}}{1 - \sigma} \dd \sigma\right) \times \lambda^{\alpha + \beta -1} R (\lambda ; A) (- A x) \dd \lambda .
\end{aligned}
\]

The coefficient () above is evaluated as $\pi^{-1}\sin\pi(\alpha+\beta)$, as may be seen by expanding $(1-\sigma)^{-1}$ into powers of $\sigma$. We have thus proved \eqref{eq:IX.11.28}. 

To prove \eqref{eq:IX.11.29}, we shall make use of $\int_{0}^{\infty} \lambda^{\alpha -1}(1 + \lambda)^{-1} \dd\lambda = \pi / \sin \alpha \pi$. Thus 
\[
(- A) ^{\alpha} x - (- A) x = \frac {\sin \alpha \pi}{\pi} \int_{0} ^{\infty} \lambda^{\alpha -1} (R (\lambda ; A) - \frac {1}{\lambda + 1}) (- A x) \dd \lambda .
\]

We split the integral into two parts, one from $0$ to $C$ and the other from $C$ to $\infty$. For a fixed $C$, the first part goes to zero as $\alpha\uparrow1$, since $R(\lambda;A)(-Ax)=x-\lambda R(\lambda;A)x$ is bounded in $\lambda>0$. The second part is, in norm, 
\[
\leqq \frac {\sin \alpha \pi}{\pi (1 - \alpha)} C ^{\alpha -1} \sup\limits _{\lambda \geqq C} \left\| (\lambda R (\lambda ; A) - \frac {\lambda}{\lambda + 1}) A x \right\|.
\]

We have $\slim\limits_{\lambda \uparrow \infty} \lambda R(\lambda; A) x = x$ by $x - \lambda R(\lambda; A) x = R(\lambda; A) (-A x)$ and \eqref{eq:IX.11.27}. Hence $\slim\limits_{\alpha \uparrow 1}$ of the second part is arbitrarily near $0$ if we take $C$ sufficiently large. This prove \eqref{eq:IX.11.29}. 

To prove \eqref{eq:IX.11.30}, we split the integral into two parts, one from $0$ to $C$ and the other from $C$ to $\infty$. Because of \eqref{eq:IX.11.27}, the second part tends to zero as $\alpha \downarrow 0$. By $R(\lambda; A)(-Ax) = x - \lambda R(\lambda; A)x$ and the hypothesis $\slim\limits_{\lambda \downarrow 0} \lambda R(\lambda; A)x = 0$, the first is, for sufficiently small $C$, arbitrarily near the value $(\alpha\pi)^{-1} \sin \alpha\pi \cdot C^{\alpha}x$, which tends to $x$ when $\alpha \downarrow 0$. This proves \eqref{eq:IX.11.30}. 

We shall prove \eqref{eq:IX.11.31}. By virtue of the representation \eqref{eq:IX.11.6}, we obtain 
\[
\begin{aligned}
(\mu I - (A _{\alpha}) _{\beta}) ^{-1}& = \int_{0} ^{\infty} \int_{0} ^{\infty} (2 \pi i) ^{-2} \left(\frac {1}{\mu - \lambda^{\beta} \cdot e ^{- i \pi \beta}} - \frac {1}{\mu - \lambda^{\beta} \cdot e ^{i \pi \beta}}\right)
\\
&\quad \times \left(\frac {1}{\lambda - \zeta^{\alpha} \cdot e ^{- i \pi \alpha}} - \frac {1}{\lambda - \zeta^{\alpha} \cdot e ^{i \pi \alpha}}\right) (\zeta I - A) ^{-1} \dd \lambda \dd \zeta .
\end{aligned}
\]

This double integral is absolutely convergent in norm, and so we may interchange the order of integration. Hence we obtain \eqref{eq:IX.11.31}, because the inner integral is 
\[
\begin{aligned}
&= (2 \pi i) ^{-2} \int_{C} \frac {- 1}{\mu - z ^{\beta}} \left(\frac {1}{z - \zeta^{\alpha} e ^{- i \pi \alpha}} - \frac {1}{z - \zeta^{\alpha} e ^{i \pi \alpha}}\right) \dd z
\\
&= (2 \pi i) ^{-1} \left(\frac {- 1}{\mu - \zeta^{\alpha \beta} e ^{- i \pi \alpha \beta}} - \frac {- 1}{\mu - \zeta^{\alpha \beta} e ^{i \pi \alpha \beta}}\right)
\end{aligned}
\]
where the path of integration $C$ runs from $\infty e^{i\pi}$ to $0$ and from $0$ to $\infty e^{-i\pi}$.
\end{proof}

An example of the fractional power. If $\alpha = 1/2$, then we have, by taking $\theta = \pi$ in \eqref{eq:IX.11.17}, 
\begin{equation}\label{eq:IX.11.32}
f _{t,1/2} (s) = \pi^{-1} \int_{0} ^{\infty} e ^{-sr} \sin \left(t r ^{1/2}\right) \dd r = \pi^{-1} \sqrt {\pi} t \left(2 ^{3} \sqrt {s}\right) ^{-3} \cdot e ^ {- t ^{2} / 4 s}.
\end{equation}

Thus, if we take the semi-group $\{T_t\}$ associated with the Gaussian kernel: 
\[
\left(T _{t} x\right) (u) = \frac {1}{2 \sqrt {\pi s}} \int_{- \infty} ^{\infty} e ^ {- (u - v) ^{2} / 4 s} x (v) \dd v, x \in C [ - \infty , \infty ],
\]
then 
\[
\begin{aligned}
&\left(\hat {T} _{t,1/2} x\right) (u)  = \int_{- \infty} ^{\infty} \left\{\int_{0} ^{\infty} x (v) \frac {t}{4 \pi s ^{2}} e ^ {- ((u - v) ^{2} + t ^{2}) / 4 s} \dd s \right\} \dd v = \frac {t}{\pi} \int_{- \infty} ^{\infty} \frac {1}{t ^{2} + (u - v) ^{2}} x (v) \dd v,
\end{aligned}
\]
that is, the semi-group $\{\hat{T}_t\}$ is associated with the Poisson kernel. In this case, the infinitesimal generator $A$ of $T_{t}$ is given by the differential operator $d^2 /\dd s^2$, while the infinitesimal generator $\hat{A}$ of $\hat{T}_t$ is given by the singular integral operator 
\[
(\hat {A} _{1/2} x) (s) = \slim _{h \downarrow 0} \frac {1}{\pi} \int_{- \infty} ^{\infty} \frac {x (s - v) - x (s)}{v ^{2} + h ^{2}} \dd v,
\]
and not by the differential operator $\dd/\dd s$. For another example, see K. \textsc{Yosida} \cite{yosida373}. 


\yossection{12}{The Convergence of Semi-groups. The Trotter-Kato Theorem}
\label{sec:IX.12}

We shall denote by $\exp(tA)$ the semi-group of class $(C_{0})$ with the infinitesimal generator $A$. Concerning the convergence of semi-groups, we have 

\begin{theorem}\label{thm:IX.12.Theorem.1} Let $X$ be a locally convex, sequentially complete complex linear space. Let $\{\exp(tA_{n})\} \subseteq L(X, X)$ be a sequence of equi-continuous semi-groups of class $(C_{0})$ such that the family of operators $\{\exp(tA_{n})\}$ is equi-continuous in $t \geqq 0$ and in $n = 1, 2, \ldots$. Thus we assume that, for any continuous semi-norm $p(x)$ on $X$, there exists a continuous semi-norm $q(x)$ on $X$ such that 
\begin{equation}\label{eq:IX.12.1}
p(\exp(tA_{n})x) \leqq q(x) \text{ for all } t \geqq 0, x \in X \text{ and } n = 1, 2, \ldots.
\end{equation}
Suppose that, for some $\lambda_0$ with $\operatorname{Re}(\lambda_0) > 0$, 
\begin{equation}
\label{eq:IX.12.2}
\begin{minipage}{0.92\linewidth}
\everymath{}\let\limits\nolimits
$\lim\limits_{n\to \infty}R(\lambda_0;A_n)x = J(\lambda_0)x$ exists for all $x\in X$ in such a way that the range $R(J(\lambda_0))$ is dense in $X.$
\end{minipage}
\end{equation}
Then $J(\lambda_0)$ is the resolvent of the infinitesimal generator $A$ of an equicontinuous semi-group $\exp(tA)$ of class $(C_0)$ in $X$ and 
\begin{equation}\label{eq:IX.12.3}
\lim\limits _{n \rightarrow \infty} \exp (t A _{n}) x = \exp (t A) x \text{ for every } x \in X.
\end{equation}
\end{theorem}

Moreover, the convergence in \eqref{eq:IX.12.3} is uniform in $t$ on every compact interval of $t$. 

For the proof, we prepare 

\begin{lemmau} Let $T_{t} = \exp(tA)$ be an equi-continuous semi-group of class $(C_0)$ in $X$. Then, for any continuous semi-norm $p(x)$ on $X$, there exists a continuous semi-norm $q(x)$ on $X$ such that 
\begin{equation}\label{eq:IX.12.4}
p\left(T_{t}x - (I - t n^{-1}A)^{-n}x\right) \leqq (2n)^{-1}t^{2}q\left(A^{2}x\right) \quad (n = 1, 2, \ldots)
\end{equation}
whenever $x \in D(A^2)$. 
\end{lemmau}


\begin{proof}
 Set $T(t, n) = (I - n^{-1}tA)^{-n}$. Then we know (Chapter~\ref{chap:IX}, \ref{sec:IX.7}) that $\{T(t, n)\}$ is equi-continuous in $t \geqq 0$ and $n = 1, 2, \ldots$. Moreover, we have (Chapter~\ref{chap:IX}, \ref{sec:IX.4}), for any $x \in D(A)$, 
\begin{equation}D _{t} T (t, n) = (I - n ^{-1} t A) ^{-n-1} A x = A (I - n ^{-1} t A) ^{-n-1} x, D _{t} T x = T _{t} A x = A T _{t} x. \notag\end{equation}

Thus, by the commutativity of $T_{t}$ and $T(t, n)$, 
\begin{equation}\label{eq:IX.12.5}
\begin{aligned} T _{t} x - T (t, n) & x = \int_{0} ^{t} [ D _{s} T (t - s, n) T _{s} x ] \dd s \\ & = \int_{0} ^{t} T (t - s, n) T _{s} \left(A x - \left(I - \frac {t - s}{n} A\right) ^{-1} A x\right) \dd s, x \in D (A). \end{aligned}
\end{equation}

Hence, if $x \in D(A^2)$, we have, by $(I - m^{-1}A)^{-1}Ax = -m(I - (I - m^{-1}A)^{-1})x$, 
\[
p (T _{t} x - T (t, n) x) \leqq \int_{0} ^{t} p [ T (t - s, n) T _{s} (I - n ^{-1} (t - s) A) ^{-1} \frac {s - t}{n} A ^{2} x ] \dd s,
\]
and so, by a continuous semi-norm $q(x)$ on $X$ which is independent of $x$ and $n$, 
\[
p \left(T _{t} x - T (t, n) x\right) \leqq (2 n) ^{-1} t ^{2} q \left(A ^{2} x\right).
\]
\end{proof}

\begin{corollaryu} For any $x \in D(A^2)$, $s > 0$ and $t \geqq 0$, 
\begin{equation}\label{eq:IX.12.6}
p \left(T _{t} x - (I - s A) ^{-[t/s]} x\right) \leqq s q _{1} (A x) + \frac {t s}{2} q \left(A ^{2} x\right), \text{ where }
\end{equation}
$q_{1}(x)$ is a continuous semi-norm on $X$ which is independent of $x, t$ and $s$, and $[t / s]$ is the largest integer $\leqq t / s$. 
\end{corollaryu}

\begin{proof}
 We have, for $t = ns$, 
\[
p \left(T _{ns} x - (I - s A) ^{-n} x\right) \leqq 2 ^{-1} s t q \left(A ^{2} x\right). \]

If $t = ns + u$ with $0 \leqq u < s$ and $n = [t / s]$, 
\[
p \left(T _{t} x - T _{ns} x\right) = p \left(\int_{ns} ^{t} T _{\sigma} ^{\prime} x \dd \sigma\right) \leqq \int_{ns} ^{t} p \left(T _{\sigma} A x\right) \dd \sigma \leqq s q _{1} (A x).
\]
\end{proof}

\begin{proof}[Proof of Theorem~\ref{thm:IX.12.Theorem.1}.]
By \eqref{eq:IX.12.1} and by \eqref{eq:IX.4.11} in Chapter~\ref{chap:IX}, \ref{sec:IX.4}, we see that $\{(\operatorname{Re}(\lambda) R(\lambda; A_n)^m)\}$ is equi-continuous in $\operatorname{Re}(\lambda) > 0$, in $n = 1, 2, \ldots$ and in $m = 0, 1, 2, \ldots$. From this and by \eqref{eq:IX.12.2}, we can prove that $J(\lambda_0) = (\lambda_0 I - A)^{-1}$ with some $A$ and
\begin{equation}
\label{eq:IX.12.7}
\begin{minipage}{0.92\linewidth}
\everymath{}\let\limits\nolimits
$\lim\limits _{n \rightarrow \infty} R (\lambda; A _{n}) x = R (\lambda; A) x$ whenever $\operatorname{Re} (\lambda) > 0,$ and the convergence is uniform with respect to $\lambda$ on every compact subset of the right half-plane $\operatorname{Re} (\lambda) > 0.$
\end{minipage}
\end{equation}

To this purpose, we observe that 
\[
R (\lambda ; A _{n}) x = \sum\limits_{m=0} ^{\infty} (\lambda_{0} - \lambda) ^{m} R (\lambda_{0}; A _{n}) ^{m+1} x \text{ (for } | \lambda - \lambda_{0} | / \operatorname{Re} (\lambda_{0}) < 1)
\]
and that the series is, by the equi-continuity of $\{(\operatorname{Re}(\lambda_0)R(\lambda_0^{\prime};A_n))^m\}$ in $n = 1,2,\ldots$ and $m = 0,1,2,\ldots$, uniformly convergent for $|\lambda -\lambda_0| / \operatorname{Re}(\lambda_0)\leqq 1 - \varepsilon$ and $n = 1,2,\ldots$. Here $\varepsilon$ is a fixed positive number. Therefore, for any $\delta >0$, there exists an $m_0$ and a continuous semi-norm $q(x)$ on $X$ such that, for $|\lambda -\lambda_0| / \operatorname{Re}(\lambda_0)\leqq 1 - \varepsilon$, 
\begin{align}
p \left(R (\lambda ; A _{n}) x - R (\lambda ; A _ {n ^{\prime}}) x\right)& \leqq \sum\limits_{m=0} ^ {m _{0}} | \lambda_{0} - \lambda | ^{m}. \notag \\
p \left(R \left(\lambda_{0}; A _{n}\right) ^ {m _{0} + 1} x - R \left(\lambda_{0}; A _ {n ^{\prime}}\right) ^ {m _{0} + 1} x\right) + 2 \delta q (x) \text{ for all } x& \in X. \notag
\end{align}

Hence, by \eqref{eq:IX.12.2}, we see that $\lim\limits_{n\to \infty}R(\lambda ;A_n)x = J(\lambda)x$ exists uniformly for $|\lambda -\lambda_0| / \operatorname{Re}(\lambda_0)\leqq 1 - \varepsilon$. In this way, extending the convergence domain of the sequence $\{R(\lambda ;A_n)\}$, we see that 
$\lim\limits_{n\to \infty}R(\lambda ;A)x = J(\lambda)x$ exists and the convergence is uniform 
on any compact set of $\lambda$ in the right half-plane $\operatorname{Re}(\lambda) > 0$. 

Thus $J(\lambda)$ is a pseudo-resolvent, because $J(\lambda)$ satisfies the resolvent equation with $R(\lambda; A_n)$ ( $n = 1, 2, \ldots$). However, by $R(J(\lambda_0))^a = X$ and the ergodic theorem for pseudo-resolvents in Chapter~\ref{chap:VIII}, \ref{sec:VIII.4}, we see that $J(\lambda)$ is the resolvent of a closed linear operator $A$ in such a way that $J(\lambda) = R(\lambda; A)$ and $D(A) = R(R(\lambda; A))$ is dense in $X$. 

Thus we see that $\exp(tA)$ is an equi-continuous semi-group of class $(C_{0})$ in $X$. We have to show that \eqref{eq:IX.12.3} is true. But we have, by \eqref{eq:IX.12.6}, 
\[
\begin{aligned} & p \left((\exp (t A _{n}) - (I - s A _{n}) ^{-[t/s]}) (I - A _{n}) ^{-2} x\right) \\ & \leqq s q _{1} (A _{n} (I - A _{n}) ^{-2} x) + 2 ^{-1} t s q (A _{n} ^{2} (I - A _{n}) ^{-2} x), \end{aligned}
\]
for any $x \in X$, $s > 0$ and $t \geqq 0$. 

The operators 
\[
\begin{aligned} A _{n} (I - A _{n}) ^{-1} & = (I - A _{n}) ^{-1} - I, A _{n} (I - A) ^{-2} = A _{n} (I - A _{n}) ^{-1} (I - A _{n}) ^{-1} \\ & \text{ and } A _{n} ^{2} (I - A _{n}) ^{-2} = (A _{n} (I - A _{n}) ^{-1}) ^{2} \end{aligned}
\]
are equi-continuous in $n = 1,2,\ldots$. On the other hand, by \eqref{eq:IX.12.7}, 
\[
\lim\limits _{n \rightarrow \infty} (I - s A _{n}) ^{-[t/s]} (I - A _{n}) ^{-2} x = (I - s A) ^{-[t/s]} (I - A) ^{-2} x
\]
uniformly in $s$ and $t$, if $s > 0$ is bounded away from $0$ and $\infty$, and if $t$ runs over a compact interval of $[0, \infty)$. Moreover, we have, by \eqref{eq:IX.12.6}, 
\[
\begin{aligned}
&\phi \left((\exp (t A) - (I - s A) ^{-[t/s]}) (I - A) ^{-2} x\right) \leqq s q _{1} (A (I - A) ^{-2} x) + 2 ^{-1} t s q (A ^{2} (I - A) ^{-2} x)
\end{aligned}
\]
for every $x \in X$, $s > 0$ and $t \geqq 0$. Thus, by taking $s > 0$ sufficiently small, we see that 
\[
\lim\limits _{n \rightarrow \infty} \exp (t A _{n}) y = \exp (t A) y \text{ for any } y \in R (1; A) ^{2} \cdot X,
\]
and the convergence is uniform in $t$ on every compact interval of $t$. $R(1;A)^{2} \cdot X$ being dense in $X$, we see, by the equi-continuity of $\exp(tA)$ and $\exp(tA_{n})$ in $t \geqq 0$ and in $n = 1, 2, \ldots$, that \eqref{eq:IX.12.3} holds. 
\end{proof}

\begin{theorem}\label{thm:IX.12.Theorem.2} Let a sequence $\{\exp (tA_n)\}$ of equi-continuous semi-groups of class $(C_0)$ in $X$ be such that $\{\exp (tA_n)\}$ is equi-continuous in $t\geqq 0$ and in $n = 1,2,\ldots$. If, for each $x\in X$, 
\[
\lim\limits _{n \rightarrow \infty} \exp (t A _{n}) x = \exp (t A) x
\]
uniformly with respect to $t$ on every compact interval of $t$, then 
$\lim\limits_{n \to \infty} R(\lambda; A_n) x = R(\lambda; A) x$ for each $x \in X$ and $\operatorname{Re}(\lambda) > 0$ and the convergence is uniform on every compact set of $\lambda$ in the right half-plane $\operatorname{Re}(\lambda) > 0$. 
\end{theorem}

\begin{proof}
 We have 
\[
R (\lambda ; A) x - R (\lambda ; A _{n}) x = \int_{0} ^{\infty} e ^{- \lambda t} (\exp (t A) - \exp (t A _{n})) x \dd t. \]

Hence, splitting the integral into two parts, one from $0$ to $C$ and the other from $C$ to $\infty$, we obtain the result. 
\end{proof}

\begin{remarku} For the case of a Banach space $X$, Theorem~\ref{thm:IX.12.Theorem.1} was first proved by H. F. \textsc{Trotter} \cite{trotter327}. In his paper, the proof that $J(\lambda)$ is the resolvent $R(\lambda; A)$ is somewhat unclear. This was pointed out by T. \textsc{Kato}. The proof given above is adapted from Kato's modification of Trotter's proof. For perturbation of semi-groups, see E. \textsc{Hille}-R. S. \textsc{Phillips} \cite{hille97}, T. \textsc{Kato} \cite{kato142} and K. \textsc{Yosida} \cite{yosida374}. 
\end{remarku}


\yossection{13}{Dual Semi-groups. Phillips' Theorem}
\label{sec:IX.13}

Let $X$ be a locally convex sequentially complete linear topological space, and $\{T_t; t \geqq 0\} \subseteq L(X, X)$ an equi-continuous semi-group of class $(C_0)$. Then the family $\{T_t^*; t \geqq 0\}$ of operators $\in L(X', X')$ where (*) denotes the dual operators in this section, satisfies the semi-group property: $T_t^* T_s^* = T_{t+s}^*$, $T_0^* = I^* =$ the identity in $X_s'$ (see Theorem~\ref{thm:VII.1.Theorem.3} in Chapter~\ref{chap:VII}, \ref{sec:VII.1}). However, it is not of class $(C_0)$ in general. For, the mapping $T_t \to T_t^*$ does not necessarily conserve the continuity in $t$ (see Proposition~\ref{prop:VII.1.Proposition.1} in Chapter~\ref{chap:VII}, \ref{sec:VII.1}). But we can show that $\{T^*\}$ is equicontinuous in $t \geqq 0$. For, we can prove 

\begin{proposition}\label{prop:IX.13.Proposition.1} If $\{S_t; t \geqq 0\} \subseteq L(X, X)$ is equi-continuous in $t \geqq 0$, then $\{S_t^*; t \geqq 0\} \subseteq L(X', X')$ is also equi-continuous in $t \geqq 0$. 
\end{proposition}

\begin{proof}
 For any bounded set $B$ of $X$, the set $\bigcup\limits_{t\geqq 0}S_t\cdot B$ is by hypothesis a bounded set of $X$. Let $U'$ and $V'$ be the polar sets of $B$ and $\bigcup\limits_{t\geqq 0}S_t\cdot B$: $U' = \{x' \in X'; \sup\limits_{b\in B}|\langle b,x'\rangle |\leqq 1\}$, $V' = \{x' \in X'; \sup\limits_{b\in B,t\geqq 0}|\langle S_t\cdot b,x'\rangle |\leqq 1\}$. 

Then (see Chapter~\ref{chap:IV}, \ref{sec:IV.7}), $U'$ and $V'$ are neighbourhoods of $0$ of $X_s'$. From 
\[
\left| \langle S _{t} \cdot b, x ^{\prime} \rangle \right| = | < b, S _{t} ^{*} x ^{\prime} \rangle | \leqq 1 (\text{ when } b \in B, x ^{\prime} \in V ^{\prime})
\]
we see that $S_{t}^{*} \cdot V' \subseteq U'$ for all $t \geqq 0$. This proves that $\{S_{t}^{*}\}$ is equicontinuous in $t \geqq 0$. 

Let $A$ be the infinitesimal generator of the semi-group $T_{t}$. Then $D(A)^{a} = X$, $R(A) \subseteq X$, and, for $\lambda > 0$, the resolvent $(\lambda I - A)^{-1} \in L(X, X)$ exists in such a way that 
\begin{equation}
\label{eq:IX.13.1}
\begin{minipage}{0.92\linewidth}
\everymath{}\let\limits\nolimits
$\{\lambda^m (\lambda I - A)^{-m}\}$ is equi-continuous in $\lambda >0$ and in $m = 0,1\ldots$
\end{minipage}
\end{equation}
We can prove (cf. Theorem~\ref{thm:VIII.6.Theorem.2} in Chapter~\ref{chap:VIII}, \ref{sec:VIII.6}) 
\end{proof}


\begin{proposition}\label{prop:IX.13.Proposition.2} For $\lambda > 0$, the resolvent $(\lambda I^{*} - A^{*})^{-1}$ exists and 
\begin{equation}\label{eq:IX.13.2}
(\lambda I ^{*} - A ^{*}) ^{-1} = ((\lambda I - A) ^{-1}) ^{*}.
\end{equation}
\end{proposition}

\begin{proof}
 We have $(\lambda I - A)^{*} = \lambda I^{*} - A^{*}$. Since $(\lambda I - A)^{-1} \in L(X, X)$, the operator $((\lambda I - A)^{-1})^{*} \in L(X', X')$ exists. We shall prove that $(\lambda I^{*} - A^{*})^{-1}$ exists and equals $((\lambda I - A)^{-1})^{*}$. Suppose there exists an $x' \in X'$ such that $(\lambda I^{*} - A^{*}) x' = 0$. Then $0 = \langle x, (\lambda I^{*} - A^{*}) x' \rangle = \langle (\lambda I - A) x, x' \rangle$ for all $x \in D(A)$. But, since $R(\lambda I - A) = X$, we must have $x' = 0$. Hence the inverse $(\lambda I^{*} - A^{*})^{-1}$ must exist. We have, for $x \in X$, $x' \in D(A^{*})$, 
\[
\langle x, x ^{\prime} \rangle = \left\langle (\lambda I - A) (\lambda I - A) ^{-1} x, x ^{\prime} \right\rangle = \left\langle (\lambda I - A) ^{-1} x, (\lambda I ^{*} - A ^{*}) x ^{\prime} \right\rangle . \]

Thus $D(((\lambda I - A)^{-1})^*) \supseteq R(\lambda I^* - A^*)$ and $((\lambda I - A)^{-1})^* \cdot (\lambda I^* - A^*) x' = x'$ for every $x' \in D(A^*)$. This proves that $((\lambda I - A)^{-1})^* \supseteq (\lambda I^* - A^*)^{-1}$. On the other hand, if $x \in D(A)$ and $x' \in D\left((\lambda I - A)^{-1})^*\right)$, then 
$\langle x, x' \rangle = \langle (\lambda I - A)^{-1} (\lambda I - A) x, x' \rangle = \langle (\lambda I - A) x, ((\lambda I - A)^{-1})^* x' \rangle$. This proves that $D(A^*) = D((\lambda I - A)^*) \supseteq R(((\lambda I - A)^{-1})^*)$ and $(\lambda I - A)^* \cdot ((\lambda I - A)^{-1})^* x' = x'$ for every $x' \in D(((\lambda I - A)^{-1})^*)$, that is, $((\lambda I - A)^{-1})^* \subseteq (\lambda I^* - A^*)^{-1}$. We have thus proved \eqref{eq:IX.13.2}. 
\end{proof}

We are now ready to prove 

\begin{theoremu} Let $X$ be a locally convex sequentially complete linear topological space such that its strong dual space $X'$ is also sequentially complete. Let $\{T_{t}\} \subseteq L(X, X)$, $t \geqq 0$, be an equi-continuous semi-group of class $(C_{0})$ with the infinitesimal generator $A$. Let us denote by $X^{+}$ the closure $D(A^{*})^{a}$ of the domain $D(A^{*})$ in the strong topology of $X'$. Let $T_{t}^{+}$ be the restriction of $T_{t}^{*}$ to $X^{+}$. Then $T_{t}^{+} \in L(X^{+}, X^{+})$ and $\{T_{t}^{+}; t \geqq 0\}$ is an equi-continuous semi-group of class $(C_{0})$ such that its infinitesimal generator $A^{+}$ is the largest restriction of $A^{*}$ with domain and range in $X^{+}$. 
\end{theoremu}

\begin{remarku} The above Theorem was proved by R. S. \textsc{Phillips} \cite{phillips252} for the special case of a $B$-space $X$. The extension given above is due to H. \textsc{Komatsu} \cite{komatsu156}. 
\end{remarku}

\begin{proof}[Proof of the Theorem.]
We have the resolvent equation $R(\lambda; A) - R(\mu; A) = (\mu - \lambda)R(\lambda; A)R(\mu; A)$ and the equi-continuity of $\{\lambda^m R(\lambda; A)^m\}$ in $\lambda > 0$ and in $m = 0, 1, 2, \ldots$. Thus, by Proposition~\ref{prop:IX.13.Proposition.1} and $2$,
\begin{equation}\label{eq:IX.13.3}
(\lambda I ^{*} - A ^{*}) ^{-1} - (\mu I ^{*} - A ^{*}) ^{-1} = (\mu - \lambda) (\lambda I ^{*} - A ^{*}) ^{-1} (\mu I ^{*} - A ^{*}) ^{-1}
\end{equation}
$\{\lambda^m (\lambda I^* - A^*)^{-m}\}$ is equi-continuous in $\lambda > 0$ and in $m = 0,1,2,\ldots$ ($4$) 

Thus, if we denote by $J(\lambda)$ the restriction to $X^{+}$ of $(\lambda I^{*} - A^{*})^{-1}$, we have 
\begin{equation}\tag{3'}\label{eq:IX.13.3'}
J (\lambda) - J (\mu) = (\mu - \lambda) J (\lambda) J (\mu),
\end{equation}
$\{\lambda^m J(\lambda)^m\}$ is equi-continuous in $\lambda > 0$ and in $m = 0, 1, \ldots$ ($4$') 

Since $D(A^{*})$ is dense in $X^{+}$ and since $(4')$ holds, we see, as in Chapter~\ref{chap:IX}, \ref{sec:IX.7}, that $\lim\limits_{\lambda \to \infty} \lambda J(\lambda) x = x$ for $x \in X^{+}$. Thus we have $R(J(\lambda))^{a} = X^{+}$ and so, by $(7')$ in Chapter~\ref{chap:VIII}, \ref{sec:VIII.4}, that $N(J(\lambda)) = \{0\}$. Thus the pseudo-resolvent $J(\lambda)$ must be the resolvent of a closed linear operator $A^{+}$ in $X^{+}$. Hence, by the sequential completeness of $X^{+}$ and $(4')$, $A^{+}$ is the infinitesimal generator of an equi-continuous semi-group of class $(C_0)$ of operators $T_t^+ \in L(X^+, X^+)$. For any $x \in X$ and $y' \in X^{+}$, we have 
\[
\langle (I - m ^{-1} t A) ^{-m} x, y ^{\prime} \rangle = \langle x, (I ^{*} - m ^{-1} t A ^{+}) ^{-m} y ^{\prime} \rangle ,
\]
and so, by the result of the preceding section, we obtain by letting $m \to \infty$ the equality $\langle T_{t}x, y' \rangle = \langle x, T_{t}^{+}y' \rangle$. Hence $T_{t}^{*}y' = T_{t}^{+}y'$, that is, $T_{t}^{+}$ is the restriction to $X^{+}$ of $T_{t}^{*}$. 

We finally show that $A^{+}$ is the largest restriction of $A^{*}$ with domain and range in $X^{+}$. It is clear, by the above derivation of the operator $A^{+}$, that $A^{+}$ is a restriction of $A^{*}$. Suppose that $x' \in D(A^{*})$ and that $x' \in X^{+}$, $A^{*}x' \in X^{+}$. Then $(\lambda I^{*} - A^{*})x' \in X^{+}$ and hence $(\lambda I^{*} - A^{+})^{-1}(\lambda I^{*} - A^{*})x' = x'$. Thus, applying $(\lambda I^{*} - A^{+})$ from the left on both sides, we obtain $A^{*}x' = A^{+}x'$. This proves that $A^{+}$ is the largest restriction of $A^{*}$ with domain as well as range in $X^{+}$. 
\end{proof}
